% Формат файла
\documentclass[12pt]{article}
\usepackage[paperheight=297mm,
   paperwidth=210mm,
   top=20mm,
   bottom=20mm,
   left=15mm,
   right=15mm]{geometry}

% Импорт преамбулы
\usepackage{../mypreamble}

% Оглавление
\title{\textbf{Алгебра}}\date{}\author{}

\begin{document}

\maketitle
\tableofcontents
\label{toc}
\newpage

\section{Многочлены}

\begin{definition}
    Многочленом от переменной $x$ над $K$ называется выражение вида: $f(x)=a_nx^n+a_{n-1}x^{n-1}+\ldots+a_1x+a_0,\, a_k\in K\text{ -- коэффициент многочлена},\, a_n\neq0$.
\end{definition}

\begin{definition}
    Наибольшее $k$ такое, что $a_k\neq0$, называется степенью многочлена $f$:
    \begin{align*}
        &k = \deg\,f\\
        &a_0\text{ -- свободный член}\\
        &a_nx^n\text{ -- старший член}\\
        &a_n\text{ -- старший коэффициент}
    \end{align*}
\end{definition}

\begin{definition}
    Два многочлена называются равными, если их коэффициенты при соответственных степенях $x$ равны.
\end{definition}

\begin{definition}
    \begin{align*}
        &f(x)=a_nx^n+a_{n-1}x^{n-1}+\ldots+a_1x+a_0\\
        &g(x)=b_nx^n+b_{n-1}x^{n-1}+\ldots+b_1x+b_0
    \end{align*}
    Суммой многочленов $f(x)$ и $g(x)$ называется: $$n(x)=(a_n+b_n)x^n+(a_{n-1}+b_{n-1})x^{n-1}+\ldots+(a_1+b_1)x+(a_0+b_0)$$ Произведением многочленов $f(x)$ и $g(x)$ называется: $$S(x)=d_{2n}x^{2n}+d_{2n-1}x^{2n-1}+\ldots+d_1x+d_0\text{, где } d_k=a_0b_k+a_1b_{k-1}+\ldots+a_{k-1}b_1+a_kb_0=\sum_{i=0}^ka_ib_{k-i}$$
\end{definition}

\begin{statement}
    Пусть $\deg\,f(x)\neq0,\,\deg\,g(x)\neq0$, тогда:
    \begin{align*}
        &1.\,\,\deg(f(x)+g(x))\leq \max {\{\deg\,f,\,\deg\,g\}}\\
        &2.\,\,f(x)\cdot g(x) \neq0\\
        &3.\,\,\deg(f(x)\cdot g(x))=\deg\,f(x)+\deg\,g(x)
    \end{align*}
\end{statement}
\begin{proof}
    \begin{align*}
    1.\,\,&\text{Пусть }\deg\,f=\deg\,g=n\\
    &f(x)+g(x)=(a_n+b_n)x^n+(a_{n-1}+b_{n-1})x^{n-1}+\ldots+(a_1+b_1)x+(a_0+b_0)\\
    &\text{Если }k>n \text{, то } a_k=0,\,b_k=0\text{, то есть }(a_k+b_k)=0\\\\
    &\text{Пусть }\deg\,f=n,\,\deg\,g=m,\,m<n\\
    &\text{Если }k>n\text{, то }a_k+b_k=0\text{, так как }b_{m+1}=b_{m+2}=\ldots=b_{n-1}=b_n=0\\
    &\text{Тогда }a_n+b_n=a_n\neq0\\\\
    2.\,\,&f(x)\cdot g(x)=\notequalto{a_nb_mx^{n+m}}{0}+\underbrace{\ldots\ldots}_{\text{степень}<n}
    \end{align*}
\end{proof}

\begin{definition}
    Многочлен $f(x)$ делится на многочлен $g(x)$, если существует такой многочлен $h(x)$, что $h(x)\cdot g(x)=f(x)$.
\end{definition}
\begin{statement}
    Всякий многочлен $f(x)\neq0$ делится на самого себя.
\end{statement}
\begin{statement}
    Если $f(x)$ делится на $g(x)$, а $g(x)$ делится на $f(x)$, то $f(x)=c\cdot g(x),\,c\in K$.
\end{statement}
\begin{definition}
    Число $x_0$ является корнем $f(x)$, если $f(x_0)=0$.
\end{definition}
\begin{theorem}[Теорема Безу]
    Остаток от деления многочлена $P(x)$ на двучлен $(x-a)$ равен $P(a)$.
\end{theorem}
\begin{proof}
    \begin{align*}
    &P(x)=(x-a)\cdot q(x)+r\\
    &P(a)=0\cdot p(x)+r=r
\end{align*}
\end{proof}
\setcounter{subsection}{4}
\begin{consequence}
    Число $a$ является корнем многочлена $P(x)$ тогда и только тогда, когда $P(x)$ делится на $(x-a)$.
\end{consequence}

\subsection{Иррациональные уравнения}

\begin{align*}
&\sqrt{f(x)}=g(x)\Longleftrightarrow 
    \begin{cases}
        f(x)=(g(x))^2\\
        g(x)\geq0
    \end{cases} & &\sqrt{f(x)}=\sqrt{g(x)}\Longleftrightarrow
    \begin{cases}
        f(x)= g(x)\\
        f(x)\geq 0
    \end{cases}
\end{align*}

\subsection{Иррациональные неравенства}

\renewcommand*{\arraystretch}{5}
$$\begin{matrix}
    \dfrac{f(x)}{g(x)}\leq 0\Longleftrightarrow \left[
\begin{gathered}
    \begin{cases}
        f(x)\geq 0\\
        g(x)<0
    \end{cases}\\
    \begin{cases}
        f(x)\leq 0\\
        g(x)>0
    \end{cases}
\end{gathered}
\right. & \sqrt{f(x)}>g(x) \Longleftrightarrow
\begin{sqcases}
    \begin{cases}
        g(x)<0\\
        f(x)\geq 0
    \end{cases}\\
    \begin{cases}
        g(x)\geq 0\\
        f(x)>(g(x))^2
    \end{cases}
\end{sqcases}\\
\sqrt{f(x)}<g(x) \Longleftrightarrow
    \begin{cases}
        g(x)\geq 0\\
        f(x) \geq 0\\
        f(x) < (g(x))^2
    \end{cases} & 
\sqrt{f(x)}>\sqrt{g(x)} \Longleftrightarrow 
    \begin{cases}
        f(x)>g(x)\\
        g(x)\geq 0
    \end{cases}
\end{matrix}$$
\renewcommand*{\arraystretch}{1}

\subsection{Неравенства с модулем}

\begin{align*}
    &|f(x)|<a, \, a>0 \Longleftrightarrow 
    \begin{cases}
        f(x)>-a\\
        f(x)<a
    \end{cases}& 
    &|f(x)|>a \Longleftrightarrow
    \begin{sqcases}
        f(x)>a\\
        f(x)<-a
    \end{sqcases}
\end{align*}

\begin{align*}
    &|f(x)|\leq|g(x)| \Longleftrightarrow (f(x)-g(x))\cdot(f(x)+g(x))\leq0\\
    &|f(x)|+|g(x)|>|f(x)+g(x)| \Longleftrightarrow f(x)\cdot g(x) < 0\\
    &|f(x)|+|g(x)|\leq|f(x)+g(x)| \Longleftrightarrow f(x)\cdot g(x) \geq 0
\end{align*}

\section{Множества}

\begin{definition}
    Множества равномощны, если между ними существует биекция.
\end{definition}
\begin{definition}
    Множества $A$ и $B$ называются равными, если $A\subseteq B,\,B\subseteq A.$
\end{definition}
\begin{definition}
    Множества, равномощные $\N$, называются счетными.
\end{definition}
\begin{definition}
    Декартовым произведением множеств $A$ и $B$ называется множество:
    $$A\times B=\{x\mid x=(a,b),\,a\in A,\,b\in B\}$$
\end{definition}
\begin{definition}
    Число $a$ называется числом кратности $k$ многочлена $f(x)$, если $f(x)$ делится на $(x-a)^k$, но не делится на $(x-a)^{k+1}$.
\end{definition}
\begin{definition}
    Множество $m$ называется ограниченным, если:
    $$\exists\, c:\,\forall x\in M\enspace |x| \leq c$$
\end{definition}
\begin{definition}
    Пусть $M \subset \R,\,M$ ограничено. Тогда наименьшая из верхних граней множества $M$ называется точной верхней гранью (супремумом):
    $$a=\sup M\Longleftrightarrow \forall x \in M:\, x \leq a,\, \forall \varepsilon>0\enspace \exists\, x \in M:\, x>a-\varepsilon$$
\end{definition}
\begin{definition}
    Пусть $M \subset \R,\,M$ ограничено. Тогда наибольшая из нижних граней множества $M$ называется точной нижней гранью (инфинумом):
    $$a=\inf M\Longleftrightarrow \forall x \in M:\, x \geq a,\, \forall \varepsilon>0\enspace \exists\, x \in M:\, x<a+\varepsilon$$
\end{definition}
\begin{definition}
    Множество $M$ называется открытым, если:
    $$\forall\,m\in M\enspace\exists\,\varepsilon>0:\,U_\varepsilon(m)\subset M$$
\end{definition}
\begin{definition}
    Точка $A$ называется предельной точкой множества $M$, если:
    $$\forall\,\varepsilon>0\enspace U_\varepsilon(A)\cap M\text{ бесконечно или } \mathring{U}_\varepsilon(A)\cap M\neq \varnothing$$
\end{definition}
\begin{definition}
    Множество называется замкнутым, если оно содержит все свои предельные точки.
\end{definition}
\begin{definition}
    Пусть $M\subset U_{A_i}$, $A_i$ -- открытое множество, тогда $U_{A_i}$ называется открытыми покрытиями.
\end{definition}
\begin{definition}
    Говорят, что почти все элементы множества $M$ удовлетворяют некоторому условию, если этому условию не удовлетворяет конечное число элементов.
\end{definition}

\section{Числовые последовательности}

\begin{definition}
    Числовой последовательностью $a_n$ называется отображение $a:\N \to \R$.
\end{definition}
\begin{definition}
    Последовательность $a_n$ называется ограниченной, если $a(\N)$ ограничено. 
\end{definition}
\begin{definition}
    Последовательность $a_n$ называется монотонно возрастающей, если:
    $$\forall\, n \in \N:\, a_{n+1}>a_n$$
\end{definition}
\begin{definition}
    Последовательность $a_n$ асимптотически больше $b_n$, если:
    $$a_n\succ b_n\Longleftrightarrow\exists\, N\in \N:\,\forall\,n>N\enspace a_n\geq b_n$$
\end{definition}
\begin{definition}
    Точка $A$ называется частичным пределом последовательности $a_n$, если:
    $$\forall\,\varepsilon>0\enspace a^{-1}(U_\varepsilon(A))\text{ бесконечно}$$
\end{definition}
\begin{theorem}
    Пусть последовательность $a_n$ ограничена, тогда у неё есть частичный предел.
\end{theorem}
\begin{proof}
    Если $\exists\, x_0:\,a^{-1}(x_0)$ бесконечно, то доказано. Если $\exists\, x_0:\,a^{-1}(x_0)$ конечно или пусто, то:
    \begin{center}
        \begin{tikzpicture}
            \draw[-{Stealth[scale = 1.5]}] (0,0) -- (10,0) node[right] {$x$};
            \coordinate (a0) at (2,0);
            \coordinate (b0) at (8,0);
            \coordinate (c0) at (5,0);
            \coordinate (c1) at (6.5,0);
            \draw (a0) node[below] {$a_0$};
            \draw (b0) node[below] {$b_0$};
            \draw (c0) node[below] {$c_0$};
            \draw (c1) node[below] {$c_1$};
            \draw [decoration={brace,mirror,raise=1.2cm},decorate] (a0) -- (b0) node [pos=0.5,anchor=north,yshift=-1.25cm] {$I_0$};
            \draw [decoration={brace,mirror,raise=0.5cm},decorate] (c0) -- (b0) node [pos=0.5,anchor=north,yshift=-0.55cm] {$I_1$}; 
            \tkzDrawPoints(a0, b0, c0, c1)
        \end{tikzpicture}
    \end{center}
    Отметим на числовой прямой $a_0=\inf a(\N)$ и $b_0=\sup a(\N)$, а также середину $a_0b_0$, то есть $c_0=\dfrac{a_0+b_0}{2}$. Разделим один из получившихся отрезков (отметим $c_0$ и $b_0$, как $a_1$ и $b_1$ соответственно) пополам, получив $c_1=\dfrac{a_1+b_1}{2}$. Данный процесс можно продолжать, получая следующую конструкцию:
    $$[a_0;\,b_0]\supset [a_1;\,b_1]\supset \ldots \supset [a_n;\,b_n]\supset \dots$$
    Теперь необходимо доказать следующее:
    $$\bigcap_{n=0}^{\infty}[a_n;\,b_n]\neq \varnothing$$
    \begin{align*}
        1.\,\,&a_0\leq a_1 \leq \ldots \leq a_n \leq \ldots\\
        2.\,\,&a(\N)\text{ ограничена сверху } b_i\text{ элементом}\\
        3.\,\,&a(\N)\text{ имеет точную верхнюю грань } M_1=\sup a(\N)\\
        &\text{и точную нижнюю грань } M_2 = \inf a(\N)\\
        4.\,\,&M_1\leq M_2\\
        5.\,\,&[M_1;\,M_2]\subset \bigcap_{n=0}^{\infty}[a_n;\,b_n]\\
        6.\,\,&\text{Пусть } M_1<M_2\text{, тогда } \exists\, n:\, b_n-a_n<M_2-M_1.\\
        &\text{Получаем противоречие, значит } M_1=M_2=M.\\
        7.\,\,&\text{Возьмем такое } n \text{, что } b_n-a_n<\varepsilon. \text{ Тогда } [a_n;\,b_n]\subset U_{\varepsilon}(M).\\
        &\text{То есть }\forall \varepsilon>0:\, a^{-1}(U_{\varepsilon}(M))\text{ бесконечно.}
    \end{align*}
\end{proof}
\begin{theorem}
    Пусть $M$ -- множество частичных пределов последовательности $a_n$, тогда:
    $$\limsup_{n\to\infty}a_n=\sup M;\enspace\liminf_{n\to\infty}a_n=\inf M$$
\end{theorem}
\begin{theorem}
    Множество частичных пределов числовой последовательности замкнуто.
\end{theorem}
\begin{definition}
    Пусть даны $a:\,\N\to\R$ и $b:\,\N\to\N$, при этом $b$ монотонна. Тогда $a\,\circ\, b:\\\N\to\R$ называется подпоследовательностью.
\end{definition}
\begin{statement}
    У всякой последовательности есть монотонная подпоследовательность.
\end{statement}

\subsection{Прогрессии}

\begin{definition}
    Арифметической прогрессией называется числовая последовательность, заданная формулой n-го члена:
    $$a_n=a_1+(n-1)\cdot d$$
\end{definition}

\begin{definition}
    Разностью арифметической прогрессии называется разность $a_{n+1}$ и $a_n$.
\end{definition}

\begin{statement}
    Пусть $(a_n)$ -- арифметическая прогрессия, тогда:
    $$a_{n+2}-a_{n+1}=a_{n+1}-a_n$$
\end{statement}

\begin{statement}
    Пусть $(a_n)$ -- арифметическая прогрессия, тогда:
    $$\forall n \in \N,\,n\geq 2 \enspace\forall k \in \N,\, k<n:\, a_n=\frac{a_{n+k}+a_{n-k}}{2}$$
\end{statement}

\begin{theorem}
    Сумма первых $n$ членов арифметической прогрессии равна:
    $$S_n=n\cdot \left(a_1+\frac{(n-1)\cdot d}{2} \right)=n \cdot \frac{a_1 + a_n}{2}$$
\end{theorem}

\begin{proof}
    \begin{align*}
        &S_n=a_1+a_2+a_3+\ldots+a_n=\\
        =&a_1+(a_1+d)+(a_1+2d)+\ldots+(a_1+(n-1)\cdot d)=\\
        =&n\cdot a_1+d\cdot \left( \frac{(n-1)\cdot n}{2} \right)=\\
        =&n\cdot \left(a_1+\frac{(n-1)\cdot d}{2} \right)=\\
        =&n \cdot \frac{a_1 + (a_1 + (n-1)\cdot d)}{2}=\\
        =&n\cdot \frac{a_1 + a_n}{2}
    \end{align*}
\end{proof}

\begin{definition}
    Геометрической прогрессией называется числовая последовательность, заданная формулой n-го члена:
    $$b_n=b_1 \cdot q^{n-1},\, b_1\neq0,\, q\neq0$$
\end{definition}

\begin{statement}
    Пусть $(b_n)$ -- геометрическая прогрессия, тогда: 
    $$\forall n\in \N,\, n\geq 2:\,b_n^2=b_{n-1}\cdot b_{n+1}$$
\end{statement}
\begin{proof}
    $$b_{n-1}\cdot b_{n+1}=b_1\cdot q^{n-2}\cdot b_1 \cdot q^n=b_1^2\cdot q^{2n-2}=(b_1\cdot q^{n-1})^2$$
\end{proof}
\begin{theorem}
    Сумма первых $n$ членов геометрической прогрессии равна:
    $$S_n=b_1\cdot \frac{1-q^n}{1-q},\, q\neq 1$$
\end{theorem}
\begin{proof}
    \begin{align*}
        &S_n=b_1+b_2+b_3+\ldots+b_n=\\
        =&b_1+b_1\cdot q+b_1\cdot q^2+\ldots+b_1\cdot q^{n-1}=\\
        =&b_1\cdot(1+q+q^2+\ldots+q^{n-1})=\\
        =&b_1\cdot \frac{1-q^n}{1-q}
    \end{align*}
\end{proof}

\subsection{Пределы последовательностей}

\begin{definition}
    Число $A$ называется пределом последовательности $a_n$, если:
    $$A=\lim_{n\to\infty}a_n=\limsup\limits_{n\to\infty}a_n=\liminf\limits_{n\to\infty}a_n$$
\end{definition}
\begin{definition}
    Число $A$ называется пределом последовательности $a_n$, если $\forall\,\varepsilon>0$ почти все элементы последовательности принадлежат $U_\varepsilon(A)$.
\end{definition}
\begin{definition}
    Число $A$ называется пределом последовательности $a_n$, если:
    $$\forall\,\varepsilon>0\enspace\exists\,n_0\in\N:\,\forall\,n>n_0,\,n\in\N\enspace |a_n-A|<\varepsilon$$
\end{definition}
\begin{theorem}
    Пусть $\lim\limits_{n\to\infty}a_n=A$; $\lim\limits_{n\to\infty}b_n=B$. Тогда:
    $$\lim\limits_{n\to\infty}(a_n+b_n)=\lim\limits_{n\to\infty}a_n+\lim\limits_{n\to\infty}b_n=A+B$$
\end{theorem}
\begin{proof}
    \begin{multline*}
        \forall\,\varepsilon>0\enspace\exists\,N_1:\,\forall\,n>N_1\enspace |a_n-A|<\frac{\varepsilon}{2};\enspace \exists\,N_2:\,\forall\,n>N_2\enspace |b_n-B|<\frac{\varepsilon}{2}\Longrightarrow\\
        \Longrightarrow\forall\,n>N=\max(N_1;N_2)\enspace |a_n+b_n-A-B|\leq |a_n-A|+|b_n-B|=\frac{\varepsilon}{2}+\frac{\varepsilon}{2}=\varepsilon
    \end{multline*}
\end{proof}
\begin{definition}
    Последовательность $a_n$ называется бесконечно малой, если $\lim\limits_{n\to\infty}a_n=0$.
\end{definition}
\begin{theorem}
    Пусть $\alpha_n$ -- бесконечно малая последовательность, тогда:
    $$\lim_{n\to\infty}=A\Longleftrightarrow a_n=A+\alpha_n$$
\end{theorem}
\begin{theorem}
    Пусть $a_n$ сходится, $\alpha_n$ -- бесконечно малая. Тогда $b_n=a_n\cdot\alpha_n$ -- бесконечно малая.
\end{theorem}
\begin{proof}
    $$\forall\,\varepsilon>0\enspace \exists\,N:\,\forall\,n>N\enspace |\alpha_n|<\frac{\varepsilon}{c}\Longrightarrow\forall\,n>N\enspace |a_n\cdot\alpha_n|<c\cdot\frac{\varepsilon}{c}=\varepsilon$$
\end{proof}
\begin{theorem}
    Пусть $\lim\limits_{n\to\infty}a_n=A$; $\lim\limits_{n\to\infty}b_n=B$. Тогда:
    $$\lim\limits_{n\to\infty}(a_n\cdot b_n)=\lim\limits_{n\to\infty}a_n\cdot\lim\limits_{n\to\infty}b_n=AB$$
\end{theorem}
\begin{proof}
    $$\lim_{n\to\infty}(a_n\cdot b_n)=\lim_{n\to\infty}((A+\alpha_n)(B+\beta_n))=\lim_{n\to\infty}(AB+A\beta_n+B\alpha_n+\alpha_n\beta_n)=AB$$
\end{proof}
\begin{definition}
    Последовательность $x_n$ называется бесконечно большой, если:
    $$\lim_{n\to\infty}x_n=+\infty\Longleftrightarrow\forall\,c\enspace\exists\, n_0:\,\forall\,n>n_0\enspace x_n>c$$
\end{definition}
\begin{theorem}[Теорема Вейерштрасса]
    Пусть $x_n$ монотонна, тогда:
    \begin{align*}
        1.\,\,&\text{Она имеет предел в }\bar{\R}=\R\cup\{-\infty;\,+\infty\}\\
        2.\,\,&\text{Если она ограничена, то она имеет предел в }\R
    \end{align*}
\end{theorem}
\begin{proof}
    Пусть $x_n$ монотонно возрастает. По определению $\forall n:\, x_{n+1}>x_n$. Пусть $x_n$ не ограничена, то есть $\nexists\,m:\,\forall n\enspace \,x_n<m$, тогда $\sup(x_n)=+\infty$, а значит $x_n\to \infty$. Пусть $\exists\,m:\,\forall n \enspace \,x_n\leq m$ и $m=\sup(x_n)$. Тогда $m=\lim\limits_{n\to\infty}x_n$. Доказательство для монотонно убывающей последовательности аналогично.
\end{proof}
\begin{theorem}[Принцип двух милиционеров]
    Пусть даны последовательности: $x_n\leq y_n\leq z_n$; $\lim\limits_{n\to\infty}x_n=\lim\limits_{n\to\infty}z_n=A$, тогда:
    $$\lim_{n\to\infty}x_n=\lim_{n\to\infty}y_n=\lim_{n\to\infty}z_n=A$$
\end{theorem}
\begin{proof}
    $$\forall\,\varepsilon>0\enspace \exists\,N:\,\forall\,n>N\enspace |x_n-A|<\varepsilon;\enspace |z_n-A|<\varepsilon\Longrightarrow|x_n-A|\leq|y_n-A|\leq|z_n-A|<\varepsilon$$
\end{proof}
\begin{theorem}
    Пусть $a_n$ монотонно не убывает и ограничена сверху, тогда $a_n$ сходится.
\end{theorem}
\begin{proof}
    В силу ограниченности $a_n$ имеет частичный предел. Пусть $B$ -- частичный предел $a_n$ и $B<\sup a_n=A$. Возьмём $\varepsilon<\frac{a-b}{3}$, тогда $\exists\,n:\,a_n>A-\varepsilon$. Значит, $a_n-B>2\varepsilon$ и $\forall\,m>n\enspace a_m-B\geq a_n-B$, то есть $\lim\limits_{n\to\infty}a_n=\sup a_n$.
\end{proof}
\begin{lemma}[Неравенство Бернулли]
    $$(1+x)^n\geq 1+xn,\,n\in \N,\,x>0$$
\end{lemma}
\begin{proof}
    Докажем по индукции: $(1+x)^1\geq 1+1\cdot x$ верно, пусть $(1+x)^n\geq 1+xn$ верно, тогда:
    $$(1+x)^{n+1}=(1+x)(1+x)^n\geq (1+x)(1+xn)=1+x+nx+x^2n>1+(n+1)x$$
\end{proof}
\begin{definition}
    $$e=\lim_{n\to \infty}\left(1+\dfrac{1}{n}\right)^n$$
\end{definition}
\begin{proof}
    Рассмотрим последовательность $E_n$:
    $$E_n=\left(\dfrac{n+1}{n}\right)\cdot e_n=\left(1+\dfrac{1}{n}\right)^{n+1}\geq 1+\dfrac{n+1}{n};\enspace \,2<2+\frac{1}{n}\leq3$$
    $$\dfrac{E_{n+1}}{E_n}=\dfrac{\left(1+\dfrac{1}{n+1}\right)^{n+2}}{\left(1+\dfrac{1}{n}\right)^{n+1}}=\dfrac{\left(\dfrac{n+2}{n+1}\right)^{n+2}}{\left(\dfrac{n+1}{n}\right)^{n+1}}=\left(\dfrac{n+2}{n+1}\right)\cdot\left(\dfrac{n^2+2n}{n^2+2n+1}\right)^{n+1}$$
    Рассмотрим последовательность, в которой поменяем местами числитель и знаменатель:
    $$\left(\dfrac{n^2+2n+1}{n^2+2n}\right)^{n+1}=\left(1+\dfrac{1}{n^2+2n}\right)^{n+1}\geq1+\dfrac{n+1}{n^2+2n}=\dfrac{n^3+3n+1}{n^2+2n}$$
    Вернувшись к изначальной последовательности, получаем:
    $$\left(\dfrac{n^2+2n}{n^2+2n+1}\right)^{n+1}<\dfrac{n^2+2n}{n^3+3n+1}$$
    Таким образом, $E_{n+1}>E_n$, а значит $E_n$ монотонно возрастает, при этом она ограничена, а значит по теореме Вейерштрасса имеет предел, то есть $e_n$ также имеет предел.
\end{proof}

\section{Математический анализ}

\subsection{Эквивалентность и группы}

\begin{definition}
    Пусть $M$ -- множество, тогда множество $R \subset \{(a,\,b)\mid a,\,b\in M\}$ упорядоченных пар элементов $M$ называется бинарным отношением на $M$. 
\end{definition}

\begin{definition}
    Бинарное отношение называется отношением эквивалентности, если оно удовлетворяет свойствам:
    \begin{align*}
        1.\,\,&\text{Рефлексивность } a\sim a\\
        2.\,\,&\text{Симметричность } a\sim b \Longleftrightarrow b\sim a\\
        3.\,\,&\text{Транзитивность } a\sim b,\, b\sim c \Longleftrightarrow a\sim c
    \end{align*}
\end{definition}
\begin{theorem}[Малая теорема Ферма]
    $\forall n\in \N,\, p\in \prob:\, n^{p-1}\equiv_p 1$
\end{theorem}
\begin{definition}
    Бинарной операцией $\times$ на множестве $M$ называется отображение из множества упорядоченных пар $M^2=\{(a,\,b)\mid a,\,b\in M\}$ в множество $M$.
\end{definition}
\begin{definition}
    Пара $G(M;\,\times)$, $M$ -- множество, $\times$ -- бинарная операция, называется группой, если выполняются свойства:
    \begin{align*}
        1.\,\,&\forall a,\,b\in M:\,(a\times b)\in M\\
        2.\,\,&\exists\, e\in M\enspace\forall a\in M:\, e\times a=a\\
        3.\,\,&\forall a\in M\enspace \exists\, a^{-1}\in M:\, a\times a^{-1}=a^{-1}\times a=e\\
        4.\,\,&\forall a,\,b,\,c \in M:\,(a\times b)\times c=a\times(b\times c)= (a\times c)\times b
    \end{align*}
\end{definition}

\subsection{Пределы функций}

\begin{definition}[Предел функции по Коши]
    Число $A$ называется пределом $f(x)$ в точке $x_0$, если $f(x)$ определена в проколотой окрестности точки $x_0$, а также:
    $$A=\lim_{x\to x_0}f(x)\Longleftrightarrow\forall\,\varepsilon>0\enspace \exists\,\delta>0:\, \forall\,x\in \mathring{U}_\delta(x_0)\enspace f(x)\in U_\varepsilon(A)$$
\end{definition}
\begin{definition}[Предел функции по Гейне]
    Число $A$ называется пределом $f(x)$ в точке $x_0$, если $f(x)$ определена в проколотой окрестности точки $x_0$, а также:
    $$A=\lim_{x\to x_0}f(x)\Longleftrightarrow\forall\,x_n\neq x_0\enspace \lim_{n\to\infty}x_n=x_0\Longrightarrow \lim_{n\to\infty}f(x_n)=A$$
\end{definition}
\begin{statement}
    Определения предела функции по Коши и по Гейне эквивалентны.
\end{statement}
\begin{proof}
    Пусть определение по Коши верно, также пусть $x_n\neq x_0$ -- некоторая последовательность, такая что $\lim\limits_{n\to\infty}x_n=x_0$, то есть $\exists\,n_0:\,\forall\,n>n_0\enspace x_n\in U_\delta(x_0)$, но тогда по Коши $f(x_n)\in U_\varepsilon(A)$.\bigskip

    Пусть определение по Гейне верно, а по Коши нет, то есть:
    $$\exists\,\varepsilon>0:\,\forall\,\delta>0\enspace\exists\,x\in\mathring{U}_\delta(x_0):\,f(x)\notin U_\varepsilon(A)$$
    Возьмём $\delta=\frac{1}{k},\,k\in\N$, тогда $\exists\,x_n\in\mathring{U}_\delta(x_0)$, причём $f(x_n)\notin U_\varepsilon(A)$, что противоречит определению по Гейне.
\end{proof}
\begin{theorem}[Принцип двух милиционеров для функций]
    Пусть $f(x)$, $g(x)$ и $h(x)$ определены в проколотой окрестности точки $x_0$, тогда:
    $$f(x)\leq g(x)\leq h(x);\enspace \lim_{x\to x_0}f(x)=\lim_{x\to x_0}h(x)=A\Longrightarrow \lim_{x\to x_0}g(x)=A$$
\end{theorem}
\begin{proof}
    Пусть $x_n\neq x_0:\,\lim\limits_{n\to\infty}x_n=x_0$. Из условия $\lim\limits_{n\to\infty}f(x_n)=\lim\limits_{n\to\infty}h(x_n)=A$. Тогда по принципу двух милиционеров для последовательностей $\lim\limits_{n\to\infty}g(x_n)=A$.
\end{proof}
\begin{definition}
    Пусть $f(x)$ определена на некотором неограниченном множестве, тогда:
    $$\lim_{x\to+\infty}f(x)=B\Longleftrightarrow \forall\,\varepsilon>0\enspace\exists\,c:\,\forall\,x>c\enspace f(x)\in U_\varepsilon(B)$$
\end{definition}
\begin{theorem}
    Если у $f(x)$ существует предел в точке $x_0$, он единственный.
\end{theorem}
\begin{proof}
    Пусть у $f(x)$, определённой в проколотой окрестности точки $x_0$, есть два предела $A_1$ и $A_2$, тогда:
    $$\forall\,\varepsilon>0\enspace\exists\,\delta_1,\,\delta_2>0:\,\forall\,x_1\in\mathring{U}_{\delta_1}(x_0),\,x_2\in\mathring{U}_{\delta_2}(x_0)\enspace f(x_1)\in U_\varepsilon(A_1),\,f(x_2)\in U_\varepsilon(A_2)$$
    Возьмём $\delta=\min(\delta_1;\,\delta_2)$, тогда $\forall\,x\in\mathring{U}_\delta(x_0)\enspace f(x)\in U_\varepsilon(A_1),\,f(x)\in U_\varepsilon(A_2)$, но $U_\varepsilon(A_1)\cap U_\varepsilon(A_2)=\varnothing$, противоречие.
\end{proof}
\begin{definition}
    Пусть $f(x)$ определена в проколотой левой (правой) полуокрестности точки $x_0$. Число $A$ называется левым (правым) пределом $f(x)$ в точке $x_0$, если:
    $$A=\lim_{x\to x_0\pm0}f(x)\Longleftrightarrow \forall\,\varepsilon>0\enspace\exists\,\delta>0:\,\forall\,x\in\mathring{U}_\delta^\pm(x_0)\enspace f(x)\in U_\varepsilon(A)$$ 
\end{definition}
\begin{statement}
    Функция $f(x)$, определённая в проколотой окрестности точки $x_0$, имеет в ней предел тогда и только тогда, когда:
    $$\lim_{x\to x_0-0}f(x)=\lim_{x\to x_0+0}f(x)=\lim_{x\to x_0}f(x)$$
\end{statement}
\begin{definition}
    Точка $x_0$ называется изолированной точкой множества $M\subset\R$, если:
    $$\exists\,\varepsilon>0:\,\mathring{U}_\varepsilon(x_0)\cap M=\varnothing$$
\end{definition}
\begin{definition}
    Пусть $f(x)$ определена на $M\subset\R$. Она называется непрерывной в точке $x_0\in M$, если $x_0$ -- изолированная точка или $\lim\limits_{x\to x_0}f(x)=f(x_0)$.
\end{definition}
\begin{definition}[Непрерывность по Гейне]
    Пусть $f(x)$ определена на $M\subset\R$. Она называется непрерывной в точке $x_0\in M$, если:
    $$\forall\,x_n\in M:\,\lim_{n\to\infty}x_n=x_0\enspace \lim_{n\to\infty}f(x_n)=f(x_0)$$
\end{definition}
\begin{definition}[Непрерывность по Коши]
    Пусть $f(x)$ определена на $M\subset\R$. Она называется непрерывной в точке $x_0\in M$, если:
    $$\forall\,\varepsilon>0\enspace\exists\,\delta>0:\,\forall\,x\in U_\delta(x_0)\cap M\enspace f(x)\in U_\varepsilon(f(x_0))$$
\end{definition}
\begin{theorem}[1-я теорема Больцано-Коши]
    Пусть функция $f(x)$ определена и непрерывна на отрезке $\left[a;b\right]$, тогда:
    $$f(a)<0,\,f(b)>0\Longrightarrow\exists\,c\in \left[a;b\right]:\,f(c)=0$$
\end{theorem}
\begin{proof}
    Рассмотрим $d=\frac{a+b}{2}$. Если $f(d)=0$, то $c=d$. Если $f(d)<0$, то $a_1=d$, $b_1=b$, аналогично если $f(d)>0$, то $a_1=a$, $b_1=d$, при этом $d_1=\frac{a_1+b_1}{2}$. Таким образом:
    $$\bigcap_{n=1}^\infty\left[a_n;b_n\right]=\{c\};\enspace \lim_{n\to\infty}a_n=\lim_{n\to\infty}b_n=c$$
    Пусть $f(c)=k>0$, тогда $\forall\,n\enspace a_n <0\Longrightarrow\forall\,\delta>0\enspace\exists\,a_n\in U_\delta(c):\,|f(c)-f(a_n)|>\frac{k}{2}$, что невозможно, так как $\lim\limits_{n\to\infty}a_n=c$, противоречие. Аналогично при $f(c)<0$.
\end{proof}
\begin{theorem}
    Пусть $f(x)$ непрерывна на $\left[a;b\right]$, тогда она ограничена на $\left[a;b\right]$.
\end{theorem}
\begin{proof}
    Пусть $f(x)$ не ограничена на $\left[a;b\right]$. Возьмём $c=\frac{a+b}{2}$. Тогда на одном из отрезков $\left[a;c\right]$ или $\left[c;b\right]$ она не ограничена. Без ограничения общности пусть это отрезок $\left[a;c\right]$. Тогда $a_1=a$, $b_1=c$, $c_1=\frac{a_1+b_1}{2}$. Аналогично в противном случае. Таким образом:
    $$\bigcap_{n=1}^\infty\left[a_n;b_n\right]=\{k\}$$
    Возьмём бесконечно большую последовательность $d_n$. Тогда $\forall\, n\enspace  \exists\,x_n\in \left[a_n;b_n\right]:\,f(x_n)>d_n$, то есть $\lim\limits_{n\to\infty}x_n=k$; из непрерывности $\lim\limits_{n\to\infty}f(x_n)=f(k)$, но $f(x_n)$ -- бесконечно большая. Таким образом, значение функции в точке равно бесконечности, противоречие.
\end{proof}
\begin{theorem}
    Пусть $f(x)$ непрерывна на $\left[a;b\right]$, тогда она достигает своего супремума на этом отрезке.
\end{theorem}
\begin{proof}
    Возьмём $c=\frac{a+b}{2}$. Тогда:
    $$\sup_{\left[a;c\right]}f(x)\leq \sup_{\left[a;b\right]}f(x);\quad \sup_{\left[c;b\right]}f(x)\leq \sup_{\left[a;b\right]}f(x)$$
    Из данных неравенств следует, что на $\left[a;c\right]$ или на $\left[c;b\right]$ супремум равен супремуму на $\left[a;b\right]$. Без ограничения общности пусть это отрезок $\left[a;c\right]$. Тогда $a_1=a$, $b_1=c$, $c_1=\frac{a_1+b_1}{2}$. Аналогично в противном случае. Таким образом:
    $$\bigcap_{n=1}^\infty\left[a_n;b_n\right]=\{k\}$$
    Возьмём бесконечно малую последовательность $\varepsilon_n$. Тогда: 
    $$\forall\, n\enspace  \exists\,x_n\in \left[a_n;b_n\right]:\,f(x_n)>\sup\limits_{\left[a_n;b_n\right]}f(x)-\varepsilon_n\Longrightarrow\lim\limits_{n\to\infty}x_n=k;\enspace \lim\limits_{n\to\infty}f(x_n)=\sup\limits_{\left[a;b\right]}f(x)$$
    Значит, из непрерывности функции $f(k)=\sup\limits_{\left[a;b\right]}f(x)$.
\end{proof}

\subsection{Производные}

\begin{definition}
    Пусть $f(x)$ определена в окрестности точки $x_0$, тогда она называется дифференцируемой в точке $x_0$, если:
    $$f(x)=f(x_0)+c(x-x_0)+(x-x_0)\alpha_{x_0}(x)$$
    Где $\alpha_{x_0}(x)$ -- бесконечно малая в $x_0$ функция. Число $c$ -- называется производной, а $c(x-x_0)$ -- дифференциалом $f(x)$ в точке $x_0$.
\end{definition}
\begin{statement}
    Если $f(x)$ дифференцируема в точке $x_0$, то она непрерывна в этой точке.
\end{statement}
\begin{proof}
    $$\lim_{x\to x_0}f(x)=\lim_{x\to x_0}\left(f(x_0)+c(x-x_0)+(x-x_0)\alpha_{x_0}(x)\right)=f(x_0)$$
\end{proof}
\begin{statement}
    Функция $f(x)$ дифференцируема в точке $x_0$ тогда и только тогда, когда:
    $$c=\lim_{x\to x_0}\frac{f(x)-f(x_0)}{x-x_0}$$
\end{statement}
\begin{proof}
    Пусть $f(x)$ дифференцируема в точке $x_0$, тогда:
    \begin{multline*}
        c(x-x_0)=f(x)-f(x_0)-(x-x_0)\alpha_{x_0}(x)\Longrightarrow\\ \Longrightarrow \lim_{x\to x_0} c=\lim_{x\to x_0}\left(\frac{f(x)-f(x_0)}{x-x_0}-\alpha_{x_0}(x)\right)\Longrightarrow c=\lim_{x\to x_0}\frac{f(x)-f(x_0)}{x-x_0}
    \end{multline*}
    Пусть существует $c$ из условия, тогда:
    $$\frac{f(x)-f(x_0)}{x-x_0}=c+\alpha_{x_0}(x)\Longrightarrow f(x)=f(x_0)+c(x-x_0)+(x-x_0)\alpha_{x_0}(x)$$
\end{proof}
\begin{definition}
    Пусть $f(x)$ дифференцируема во всех точках множества $M$. Прозводной $f'(x)$ называется такая функция, что её значение равно производной $f(x)$ в каждой точке.
\end{definition}
\begin{theorem}[Ферма]
    Пусть $f(x)$ дифференцируема в точке $x_0$ и $x_0$ является её точкой экстремума, тогда $f'(x_0)=0$.
\end{theorem}
\begin{proof}
    Пусть $x_0$ -- точка максимума, тогда рассмотрим знаки производной в окрестности $x_0$:
    $$f'(x_0)=\lim_{x\to x_0}\frac{f(x)-f(x_0)}{x-x_0};\enspace \forall\,x\in(x_0;x_0+\delta)\enspace f'(x_0)<0;\enspace \forall\,x\in(x_0-\delta;x_0)\enspace f'(x_0)>0$$
    Отсюда $f'(x_0)=0$. Доказательство для точки минимума аналогично.
\end{proof}
\begin{theorem}[Ролля]
    Пусть $f(x)$ определена и непрерывна на $\left[a;b\right]$ и дифференцируема на $(a;b)$, а также $f(a)=f(b)$. Тогда:
    $$\exists\,c\in(a;b):\,f'(c)=0$$
\end{theorem}
\begin{proof}
    При $f(x)=\const$ на $\left[a;b\right]$ очевидно. Иначе:
    $$\exists\,c\in(a;b):\,f(c)=\sup\limits_{\left[a;b\right]}f(x)\text{ или }f(c)=\inf\limits_{\left[a;b\right]}f(x)$$
    То есть $c$ -- точка максимума или минимума, а значит $f'(c)=0$ по теореме Ферма.
\end{proof}
\begin{theorem}[Лагранжа]
    Пусть $f(x)$ определена и непрерывна на $\left[a;b\right]$ и дифференцируема на $(a;b)$, тогда:
    $$\exists\,c\in(a;b):\,f'(c)=\frac{f(b)-f(a)}{b-a}$$
\end{theorem}
\begin{proof}
    Если $f(a)=f(b)$, то по теореме Ролля. Рассмотрим функцию:
    $$g(x)=f(x)-\frac{f(b)-f(a)}{b-a}(x-a)$$
    Тогда $g(a)=g(b)=f(a)$, отсуда по теореме Ролля:
    $$\exists\,c\in(a;b):\,g'(c)=f'(c)-\frac{f(b)-f(a)}{b-a}=0$$
\end{proof}
\begin{theorem}
    Пусть $y=f(x)$ и $x=g(y)$ -- обратные функции, тогда:
    $$g'(y)=\frac{1}{f'(x)}$$
\end{theorem}
\begin{statement}
    Пусть $f(x)=x^n,\,n\in\N$, тогда $f'(x)=nx^{n-1}$.
\end{statement}
\begin{proof}
    $$\lim_{x\to x_0}\frac{x^n-x_0^n}{x-x_0}=\lim_{x\to x_0}(x^{n-1}+x^{n-2}x_0+\ldots+xx_0^{n-2}+x_0^{n-1})=nx_0^{n-1}$$
\end{proof}
\begin{statement}
    Пусть $f(x)=x^{-n},\,n\in\N$, тогда $f'(x)=-nx^{-n-1}$.
\end{statement}
\begin{proof}
    $$\lim_{x\to x_0}\frac{\frac{1}{x^n}-\frac{1}{x_0^n}}{x-x_0}=\lim_{x\to x_0}\frac{x_0^n-x^n}{x^nx_0^n(x-x_0)}=\frac{-nx_0^{n-1}}{x_0^{2n}}=-nx_0^{-n-1}$$
\end{proof}
\begin{statement}
    Пусть $f(x)=\sin x$, тогда $f'(x)=\cos x$.
\end{statement}
\begin{proof}
    $$\lim_{x\to x_0}\frac{\sin x-\sin x_0}{x-x_0}=\lim_{x\to x_0}\frac{2\cos\frac{x+x_0}{2}\sin\frac{x-x_0}{2}}{x-x_0}=\lim_{x\to x_0}\left(\cos\tfrac{x+x_0}{2}\cdot\frac{\sin\frac{x-x_0}{2}}{\frac{x-x_0}{2}}\right)=\cos x_0$$
\end{proof}
\begin{statement}
    Пусть $f(x)=\cos x$, тогда $f'(x)=-\sin x$.
\end{statement}
\begin{proof}
    $$\lim_{x\to x_0}\frac{\cos x-\cos x_0}{x-x_0}=\lim_{x\to x_0}\frac{-2\sin\frac{x+x_0}{2}\sin\frac{x-x_0}{2}}{x-x_0}=\lim_{x\to x_0}\left(-\sin\tfrac{x+x_0}{2}\cdot\frac{\sin\frac{x-x_0}{2}}{\frac{x-x_0}{2}}\right)=-\sin x_0$$
\end{proof}
\begin{statement}
    Пусть $f(x)=\tg x$, тогда $f'(x)=\dfrac{1}{\cos^2x}$.
\end{statement}
\begin{proof}
    $$\lim_{x\to x_0}\frac{\tg x-\tg x_0}{x-x_0}=\lim_{x\to x_0}\frac{\sin(x-x_0)}{\cos x\cos x_0 (x-x_0)}=\frac{1}{\cos^2x_0}$$
\end{proof}
\begin{statement}
    Пусть $f(x)=\ctg x$, тогда $f'(x)=-\dfrac{1}{\sin^2x}$.
\end{statement}
\begin{proof}
    $$\lim_{x\to x_0}\frac{\ctg x-\ctg x_0}{x-x_0}=\lim_{x\to x_0}\frac{-\sin(x-x_0)}{\sin x\sin x_0 (x-x_0)}=-\frac{1}{\sin^2x_0}$$
\end{proof}
\begin{statement}
    Пусть $f(x)=\sqrt{x}$, тогда $f'(x)=\dfrac{1}{2\sqrt{x}}$.
\end{statement}
\begin{proof}
    $$\lim_{x\to x_0}\frac{\sqrt{x}-\sqrt{x_0}}{x-x_0}=\lim_{x\to x_0}\frac{1}{\sqrt{x}+\sqrt{x_0}}=\frac{1}{2\sqrt{x_0}}$$
\end{proof}
\begin{statement}
    Пусть $f(x)=e^x$, тогда $f'(x)=e^x$.
\end{statement}
\begin{proof}
    Используем определение числа $e$:
    $$e=\lim_{x\to\infty}\left(1+\frac{1}{x}\right)^x=\lim_{x\to0}\left(1+x\right)^\frac{1}{x}$$
    $$\lim_{x\to x_0}\frac{e^x-e^{x_0}}{x-x_0}=\lim_{\Delta x\to0}\frac{e^{x_0+\Delta x}-e^{x_0}}{\Delta x}=e^{x_0}\lim_{\Delta x\to0}\frac{e^{\Delta x} - 1}{\Delta x}=e^{x_0}\lim_{\Delta x\to0}\frac{\left((1+\Delta x)^{\frac{1}{\Delta x}}\right)^{\Delta x}-1}{\Delta x}=e^{x_0}$$
\end{proof}
\begin{statement}
    Пусть $f(x)=a^x$, тогда $f'(x)=a^x\ln a$.
\end{statement}
\begin{proof}
    $$\frac{da^x}{dx}=\frac{de^{x\ln a}}{dx}=a^x\ln a$$
\end{proof}
\begin{statement}
    Пусть $f(x)=\ln x$, тогда $f'(x)=\dfrac{1}{x}$.
\end{statement}
\begin{proof}
    Пусть $y=\ln x$, $x=e^y$, тогда:
    $$\frac{d\ln x}{dx}=\frac{1}{\dfrac{de^y}{dy}}=\frac{1}{e^y}=\frac{1}{e^{\ln x}}=\frac{1}{x}$$
\end{proof}
\begin{statement}
    Пусть $f(x)=\log_a x$, тогда $f'(x)=\dfrac{1}{x\ln a}$.
\end{statement}
\begin{proof}
    $$\frac{d\log_a x}{dx}=\frac{d}{dx}\frac{\ln x}{\ln a}=\frac{1}{x\ln a}$$
\end{proof}
\begin{statement}
    Пусть $f(x)=\arcsin x$, тогда $f'(x)=\dfrac{1}{\sqrt{1-x^2}}$.
\end{statement}
\begin{proof}
    Пусть $y=\arcsin x$, $x=\sin y$, тогда:
    $$\frac{d\arcsin x}{dx}=\frac{1}{\dfrac{d\sin y}{dy}}=\frac{1}{\cos y}=\frac{1}{\cos(\arcsin x)}=\frac{1}{\sqrt{1-x^2}}$$
\end{proof}
\begin{statement}
    Пусть $f(x)=\arccos x$, тогда $f'(x)=-\dfrac{1}{\sqrt{1-x^2}}$.
\end{statement}
\begin{proof}
    Пусть $y=\arccos x$, $x=\cos y$, тогда:
    $$\frac{d\arccos x}{dx}=\frac{1}{\dfrac{d\cos y}{dy}}=-\frac{1}{\sin y}=-\frac{1}{\sin(\arccos x)}=-\frac{1}{\sqrt{1-x^2}}$$
\end{proof}
\begin{statement}
    Пусть $f(x)=\arctg x$, тогда $f'(x)=\dfrac{1}{1+x^2}$.
\end{statement}
\begin{proof}
    Пусть $y=\arctg x$, $x=\tg y$, тогда:
    $$\frac{d\arctg x}{dx}=\frac{1}{\dfrac{d\tg y}{dy}}=\cos^2y=\cos^2(\arctg x)=\frac{1}{1+x^2}$$
\end{proof}
\begin{statement}
    Пусть $f(x)=\arcctg x$, тогда $f'(x)=-\dfrac{1}{1+x^2}$.
\end{statement}
\begin{proof}
    Пусть $y=\arcctg x$, $x=\ctg y$, тогда:
    $$\frac{d\arcctg x}{dx}=\frac{1}{\dfrac{d\ctg y}{dy}}=-\sin^2y=-\sin^2(\arcctg x)=-\frac{1}{1+x^2}$$
\end{proof}

\subsection{Первообразная функция}

\begin{definition}
    Пусть $f(x)$ -- функция, $F(x)$ называется её первообразной, если $F'(x)=f(x)$.
\end{definition}
\begin{definition}
    Пусть $f(x)$ -- функция, тогда $\{F(x)\mid F'(x)=f(x)\}$ -- неопределённый интеграл.
\end{definition}

\section{Тригонометрия}

\begin{definition}
    Синусом угла в прямоугольном треугольнике называется отношение его противолежащего катета к гипотенузе.
\end{definition}
\begin{definition}
    Косинусом угла в прямоугольном треугольнике называется отношение его прилежащего катета к гипотенузе.
\end{definition}
\begin{definition}
    Тангенсом угла в прямоугольном треугольнике называется отношение его противолежащего катета к прилежащему или синуса этого угла к его косинусу.
\end{definition}
\begin{definition}
    Котангенсом угла в прямоугольном треугольнике называется отношение его прилежащего катета к противолежащему или косинуса этого угла к его синусу.
\end{definition}
\begin{statement}\label{Основное тригонометрическое тождество}
    $$\sin^2 \alpha + \cos^2 \alpha=1$$
\end{statement}
\begin{proof}
    По теореме Пифагора в прямоугольном треугольнике с катетами $a$ и $b$ и гипотенузой $c$:
    $$\sin^2 \alpha + \cos^2 \alpha=\left(\dfrac{a}{c}\right)^2+\left(\dfrac{b}{c}\right)^2=\dfrac{a^2 + b^2}{c^2}=1$$
\end{proof}
\begin{statement}
    $$\tg^2\alpha + 1=\dfrac{1}{\cos^2 \alpha}$$
\end{statement}
\begin{proof}
    Если выражение из утверждения \ref{Основное тригонометрическое тождество} разделить на $\cos^2 \alpha$, то получим данное выражение.
\end{proof}
\begin{statement}
    $$\ctg^2\alpha + 1=\dfrac{1}{\sin^2 \alpha}$$
\end{statement}
\begin{proof}
    Если выражение из утверждения \ref{Основное тригонометрическое тождество} разделить на $\sin^2 \alpha$, то получим данное выражение.
\end{proof}
\begin{statement}
    $$\cos(\alpha - \beta)=\cos\alpha\cdot\cos\beta+\sin\alpha\cdot\sin\beta$$
\end{statement}
    \begin{proof}
    \hfill
    \begin{center}
        \begin{tikzpicture}[scale=.75]
            \coordinate (O) at (0,0);
            \coordinate (M) at (150:3);
            \coordinate (N) at (35:3);
            \draw (O) -- (M) -- (N) -- cycle;
            \draw (O) circle (3);
            \draw (O) node[below left] {$O$};
            \draw (M) node[above left] {$M$};
            \draw (N) node[right] {$N$};
            \coordinate (A) at (0:3);
            \tkzLabelAngle(A,O,M){$\alpha$}
            \tkzMarkAngle[size=.4cm](A,O,M)
            \tkzLabelAngle(A,O,N){$\beta$}
            \tkzMarkAngle[size=.55cm](A,O,N)
            \tkzMarkAngle[size=.6cm](A,O,N)
            \begin{scope}[rotate=35]
                \draw[-{Stealth[scale = 1.5]}, dashed] (-4,0) -- (5, 0) node[below right] {$\cos \tilde{x}$};
                \draw[-{Stealth[scale = 1.5]}, dashed] (0,-4) -- (0,5) node[left] {$\sin \tilde{x}$};
            \end{scope}
            \draw[-{Stealth[scale = 1.5]}] (-4,0) -- (5, 0) node[below] {$\cos x$};
            \draw[-{Stealth[scale = 1.5]}] (0,-4) -- (0,5) node[left] {$\sin x$};
            \tkzDrawPoints(O, M, N)
        \end{tikzpicture}
    \end{center}
    По теореме Пифагора: 
    $$MN=\sqrt{(\cos\alpha-\cos \beta)^2+(\sin\alpha - \sin \beta)^2}=\sqrt{2\cdot(1-\cos\alpha\cdot\cos\beta - \sin\alpha\cdot\sin\beta)^2}$$
    Введем оси координат, повернутые относительно начальной на угол $\beta$, тогда:
    $$MN=\sqrt{(\cos(\alpha-\beta)-1)^2+\sin^2(\alpha - \beta)}=\sqrt{2\cdot(1-\cos(\alpha - \beta))}$$
    То есть $\sqrt{2\cdot(1-\cos\alpha\cdot\cos\beta - \sin\alpha\cdot\sin\beta)^2}=\sqrt{2\cdot(1-\cos(\alpha - \beta))}$, а значит $\cos(\alpha - \beta)=\cos\alpha\cdot\cos\beta+\sin\alpha\cdot\sin\beta$.
    \end{proof}
    \begin{statement}
        $$\cos(\alpha+\beta)=\cos\alpha\cdot\cos\beta-\sin\alpha\cdot\sin\beta$$
    \end{statement}
    \begin{proof}
        $$\cos(\alpha+\beta)=\cos(\alpha-(-\beta))=\cos\alpha\cdot\cos(-\beta)+\sin\alpha\cdot\sin(-\beta)=\cos\alpha\cdot\cos\beta-\sin\alpha\cdot\sin\beta$$
    \end{proof}
    \begin{statement}
        $$\cos\left(\dfrac{\pi}{2}-\alpha\right)=\sin \alpha$$
    \end{statement}
    \begin{proof}
        $$\cos\left(\dfrac{\pi}{2}-\alpha\right)=\cos\dfrac{\pi}{2}\cdot\cos\alpha+\sin\dfrac{\pi}{2}\cdot\sin\alpha=\sin\alpha$$
    \end{proof}
    \begin{statement}
        $$\sin(\alpha+\beta)=\sin\alpha\cdot\cos\beta+\cos\alpha\cdot\sin\beta$$
    \end{statement}
    \begin{proof}
        $$\sin(\alpha+\beta)=\cos\left(\left(\dfrac{\pi}{2}-\alpha\right)-\beta\right)=\cos\left(\dfrac{\pi}{2}-\alpha\right)\cdot \cos \beta+\sin\left(\dfrac{\pi}{2}-\alpha\right)\cdot \sin \beta=\sin\alpha\cdot\cos\beta+\cos\alpha\cdot\sin\beta$$
    \end{proof}
    \begin{statement}
        $$\sin(\alpha-\beta)=\sin\alpha\cdot\cos\beta-\cos\alpha\cdot\sin\beta$$
    \end{statement}
    \begin{proof}
        $$\sin(\alpha-\beta)=\sin(\alpha+(-\beta))=\sin\alpha\cdot\cos(-\beta)+\cos\alpha\cdot\sin(-\beta)=\sin\alpha\cdot\cos\beta-\cos\alpha\cdot\sin\beta$$
    \end{proof}
    \begin{statement}
        $$\tg(\alpha+\beta)=\dfrac{\tg\alpha+\tg\beta}{1-\tg\alpha\cdot\tg\beta}$$
    \end{statement}
    \begin{proof}
        $$\tg(\alpha+\beta)=\dfrac{\sin(\alpha+\beta)}{\cos(\alpha+\beta)}=\dfrac{\sin\alpha\cdot\cos\beta+\cos\alpha\cdot\sin\beta}{\cos\alpha\cdot\cos\beta-\sin\alpha\cdot\sin\beta}$$
        Тогда если разделить все на $\cos\alpha\cdot\cos\beta$, получим:
        $$\dfrac{\dfrac{\sin\alpha\cdot\cos\beta}{\cos\alpha\cdot\cos\beta}+\dfrac{\cos\alpha\cdot\sin\beta}{\cos\alpha\cdot\cos\beta}}{\dfrac{\cos\alpha\cdot\cos\beta}{\cos\alpha\cdot\cos\beta}-\dfrac{\sin\alpha\cdot\sin\beta}{\cos\alpha\cdot\cos\beta}}=\dfrac{\tg\alpha+\tg\beta}{1-\tg\alpha\cdot\tg\beta}$$
    \end{proof}
    \begin{statement}
        $$\tg(\alpha-\beta)=\dfrac{\tg\alpha-\tg\beta}{1+\tg\alpha\cdot\tg\beta}$$
    \end{statement}
    \begin{proof}
        $$\tg(\alpha-\beta)=\dfrac{\sin(\alpha-\beta)}{\cos(\alpha-\beta)}=\dfrac{\sin\alpha\cdot\cos\beta-\cos\alpha\cdot\sin\beta}{\cos\alpha\cdot\cos\beta+\sin\alpha\cdot\sin\beta}$$
        Тогда если разделить все на $\cos\alpha\cdot\cos\beta$, получим:
        $$\dfrac{\dfrac{\sin\alpha\cdot\cos\beta}{\cos\alpha\cdot\cos\beta}-\dfrac{\cos\alpha\cdot\sin\beta}{\cos\alpha\cdot\cos\beta}}{\dfrac{\cos\alpha\cdot\cos\beta}{\cos\alpha\cdot\cos\beta}+\dfrac{\sin\alpha\cdot\sin\beta}{\cos\alpha\cdot\cos\beta}}=\dfrac{\tg\alpha-\tg\beta}{1+\tg\alpha\cdot\tg\beta}$$
    \end{proof}
    \begin{statement}
        $$\ctg(\alpha+\beta)=\dfrac{\ctg\alpha\cdot\ctg\beta-1}{\ctg\alpha+\ctg\beta}$$
    \end{statement}
    \begin{proof}
        $$\ctg(\alpha+\beta)=\dfrac{\cos(\alpha+\beta)}{\sin(\alpha+\beta)}=\dfrac{\cos\alpha\cdot\cos\beta-\sin\alpha\cdot\sin\beta}{\sin\alpha\cdot\cos\beta+\cos\alpha\cdot\sin\beta}$$
        Тогда если разделить все на $\sin\alpha\cdot\sin\beta$, получим:
        $$\dfrac{\dfrac{\cos\alpha\cdot\cos\beta}{\sin\alpha\cdot\sin\beta}-\dfrac{\sin\alpha\cdot\sin\beta}{\sin\alpha\cdot\sin\beta}}{\dfrac{\sin\alpha\cdot\cos\beta}{\sin\alpha\cdot\sin\beta}+\dfrac{\cos\alpha\cdot\sin\beta}{\sin\alpha\cdot\sin\beta}}=\dfrac{\ctg\alpha\cdot\ctg\beta-1}{\ctg\beta+\ctg\alpha}$$
    \end{proof}
    \begin{statement}
        $$\ctg(\alpha-\beta)=-\dfrac{\ctg\alpha\cdot\ctg\beta+1}{\ctg\alpha-\ctg\beta}$$
    \end{statement}
    \begin{proof}
        $$\ctg(\alpha-\beta)=\dfrac{\cos(\alpha-\beta)}{\sin(\alpha-\beta)}=\dfrac{\cos\alpha\cdot\cos\beta+\sin\alpha\cdot\sin\beta}{\sin\alpha\cdot\cos\beta-\cos\alpha\cdot\sin\beta}$$
        Тогда если разделить все на $\sin\alpha\cdot\sin\beta$, получим:
        $$\dfrac{\dfrac{\cos\alpha\cdot\cos\beta}{\sin\alpha\cdot\sin\beta}+\dfrac{\sin\alpha\cdot\sin\beta}{\sin\alpha\cdot\sin\beta}}{\dfrac{\sin\alpha\cdot\cos\beta}{\sin\alpha\cdot\sin\beta}-\dfrac{\cos\alpha\cdot\sin\beta}{\sin\alpha\cdot\sin\beta}}=\dfrac{\ctg\alpha\cdot\ctg\beta+1}{\ctg\beta-\ctg\alpha}=-\dfrac{\ctg\alpha\cdot\ctg\beta+1}{\ctg\alpha-\ctg\beta}$$
    \end{proof}
    \begin{statement}
        $$\cos (2\alpha)=\cos^2\alpha-\sin^2\alpha=2\cos^2\alpha-1=1-2\sin^2\alpha$$
    \end{statement}
    \begin{proof}
        $$\cos (2\alpha)=\cos(\alpha + \alpha)=\cos^2\alpha-\sin^2\alpha=2\cos^2\alpha-1=1-2\sin^2\alpha$$
    \end{proof}
    \begin{statement}
        $$\sin(2\alpha)=2 \sin \alpha \cdot \cos \alpha$$
    \end{statement}
    \begin{proof}
        $$\sin(2\alpha)=\sin(\alpha+\alpha)=2 \sin \alpha \cdot \cos \alpha$$
    \end{proof}
    \begin{statement}
        $$\tg(2\alpha)=\dfrac{2\tg\alpha}{1-\tg^2\alpha}$$
    \end{statement}
    \begin{proof}
        $$\tg(2\alpha)=\tg(\alpha+\alpha)=\dfrac{\tg\alpha+\tg\alpha}{1-\tg\alpha\cdot\tg\alpha}=\dfrac{2\tg\alpha}{1-\tg^2\alpha}$$
    \end{proof}
    \begin{statement}
        $$\ctg(2\alpha)=\dfrac{\ctg^2\alpha-1}{2\ctg\alpha}$$
    \end{statement}
    \begin{proof}
        $$\ctg(2\alpha)=\dfrac{\ctg\alpha\cdot\ctg\alpha-1}{\ctg\alpha+\ctg\alpha}=\dfrac{\ctg^2\alpha-1}{2\ctg\alpha}$$
    \end{proof}
    \begin{statement}
        $$\cos(3\alpha)=4\cos^3\alpha-3\cos\alpha$$
    \end{statement}
    \begin{proof}
        \begin{align*}
            &\cos(3\alpha)=\\
            =&\cos(2\alpha+\alpha)=\\
            =&\cos(2\alpha)\cdot\cos\alpha-\sin(2\alpha)\cdot\sin\alpha=\\
            =&(2\cos^2\alpha-1)\cdot\cos\alpha-2\sin\alpha\cdot\cos\alpha\cdot\sin\alpha=\\
            =&2\cos^3\alpha-\cos\alpha\cdot(1+2\sin^2\alpha)=\\
            =&2\cos^3\alpha-\cos\alpha\cdot(3-2\cos^2\alpha)=\\
            =&4\cos^3\alpha-3\cos\alpha
        \end{align*}
    \end{proof}
    \begin{statement}
        $$\sin(3\alpha)=3\sin\alpha-4\sin^3\alpha$$
    \end{statement}
    \begin{proof}
        \begin{align*}
            &\sin(3\alpha)=\\
            =&\sin(2\alpha+\alpha)=\\
            =&\sin(2\alpha)\cdot\cos\alpha+\cos(2\alpha)\cdot\sin\alpha=\\
            =&2\sin\alpha\cdot\cos\alpha\cdot\cos\alpha+(1-2\sin^2\alpha)\cdot\sin\alpha=\\
            =&\sin\alpha\cdot(2\cos^2\alpha+1)-2\sin^3\alpha=\\
            =&\sin\alpha\cdot(3-2\sin^2\alpha)-2\sin^3\alpha=\\
            =&3\sin\alpha-4\sin^3\alpha
        \end{align*}
    \end{proof}
    \begin{statement}
        $$\cos^2\alpha=\dfrac{1+\cos(2\alpha)}{2}$$
    \end{statement}
    \begin{proof}
        $$\cos (2\alpha)=2\cos^2\alpha-1\Longleftrightarrow \cos^2\alpha=\dfrac{1+\cos(2\alpha)}{2}$$
    \end{proof}
    \begin{statement}
        $$\sin^2\alpha=\dfrac{1-\cos(2\alpha)}{2}$$
    \end{statement}
    \begin{proof}
        $$\sin^2\alpha=1-\cos^2\alpha=1-\dfrac{1+\cos(2\alpha)}{2}=\dfrac{1-\cos(2\alpha)}{2}$$
    \end{proof}
    \begin{statement}
        $$\cos\alpha+\cos\beta=2\cos\left(\dfrac{\alpha+\beta}{2}\right)\cdot\cos\left(\dfrac{\alpha-\beta}{2}\right)$$
    \end{statement}
    \begin{proof}
        \begin{align*}
            &\cos\alpha+\cos\beta=\\
            =&\cos\left(\dfrac{\alpha+\beta}{2}+\dfrac{\alpha-\beta}{2}\right)+\cos\left(\dfrac{\alpha+\beta}{2}-\dfrac{\alpha-\beta}{2}\right)=\\
            =&\cos\left(\dfrac{\alpha+\beta}{2}\right)\cdot\cos\left(\dfrac{\alpha-\beta}{2}\right)-\sin\left(\dfrac{\alpha+\beta}{2}\right)\cdot\sin\left(\dfrac{\alpha-\beta}{2}\right)+\\
            +&\cos\left(\dfrac{\alpha+\beta}{2}\right)\cdot\cos\left(\dfrac{\alpha-\beta}{2}\right)+\sin\left(\dfrac{\alpha+\beta}{2}\right)\cdot\sin\left(\dfrac{\alpha-\beta}{2}\right)=\\
            =&2\cos\left(\dfrac{\alpha+\beta}{2}\right)\cdot\cos\left(\dfrac{\alpha-\beta}{2}\right)
        \end{align*}
    \end{proof}
    \begin{statement}
        $$\cos\alpha-\cos\beta=-2\sin\left(\dfrac{\alpha+\beta}{2}\right)\cdot\sin\left(\dfrac{\alpha-\beta}{2}\right)$$
    \end{statement}
    \begin{proof}
        \begin{align*}
            &\cos\alpha+\cos\beta=\\
            =&\cos\left(\dfrac{\alpha+\beta}{2}+\dfrac{\alpha-\beta}{2}\right)-\cos\left(\dfrac{\alpha+\beta}{2}-\dfrac{\alpha-\beta}{2}\right)=\\
            =&\cos\left(\dfrac{\alpha+\beta}{2}\right)\cdot\cos\left(\dfrac{\alpha-\beta}{2}\right)-\sin\left(\dfrac{\alpha+\beta}{2}\right)\cdot\sin\left(\dfrac{\alpha-\beta}{2}\right)-\\
            -&\cos\left(\dfrac{\alpha+\beta}{2}\right)\cdot\cos\left(\dfrac{\alpha-\beta}{2}\right)-\sin\left(\dfrac{\alpha+\beta}{2}\right)\cdot\sin\left(\dfrac{\alpha-\beta}{2}\right)=\\
            =&-2\sin\left(\dfrac{\alpha+\beta}{2}\right)\cdot\sin\left(\dfrac{\alpha-\beta}{2}\right)
        \end{align*}
    \end{proof}
    \begin{statement}
        $$\sin\alpha+\sin\beta=2\sin\left(\dfrac{\alpha+\beta}{2}\right)\cdot\cos\left(\dfrac{\alpha-\beta}{2}\right)$$
    \end{statement}
    \begin{proof}
        \begin{align*}
            &\sin\alpha+\sin\beta=\\
            =&\sin\left(\dfrac{\alpha+\beta}{2}+\dfrac{\alpha-\beta}{2}\right)+\sin\left(\dfrac{\alpha+\beta}{2}-\dfrac{\alpha-\beta}{2}\right)=\\
            =&\sin\left(\dfrac{\alpha+\beta}{2}\right)\cdot\cos\left(\dfrac{\alpha-\beta}{2}\right)+\cos\left(\dfrac{\alpha+\beta}{2}\right)\cdot\sin\left(\dfrac{\alpha-\beta}{2}\right)+\\
            +&\sin\left(\dfrac{\alpha+\beta}{2}\right)\cdot\cos\left(\dfrac{\alpha-\beta}{2}\right)-\cos\left(\dfrac{\alpha+\beta}{2}\right)\cdot\sin\left(\dfrac{\alpha-\beta}{2}\right)=\\
            =&2\sin\left(\dfrac{\alpha+\beta}{2}\right)\cdot\cos\left(\dfrac{\alpha-\beta}{2}\right)
        \end{align*}
    \end{proof}
    \begin{statement}
        $$\sin\alpha-\sin\beta=2\cos\left(\dfrac{\alpha+\beta}{2}\right)\cdot\sin\left(\dfrac{\alpha-\beta}{2}\right)$$
    \end{statement}
    \begin{proof}
        \begin{align*}
            &\sin\alpha-\sin\beta=\\
            =&\sin\left(\dfrac{\alpha+\beta}{2}+\dfrac{\alpha-\beta}{2}\right)-\sin\left(\dfrac{\alpha+\beta}{2}-\dfrac{\alpha-\beta}{2}\right)=\\
            =&\sin\left(\dfrac{\alpha+\beta}{2}\right)\cdot\cos\left(\dfrac{\alpha-\beta}{2}\right)+\cos\left(\dfrac{\alpha+\beta}{2}\right)\cdot\sin\left(\dfrac{\alpha-\beta}{2}\right)+\\
            -&\sin\left(\dfrac{\alpha+\beta}{2}\right)\cdot\cos\left(\dfrac{\alpha-\beta}{2}\right)+\cos\left(\dfrac{\alpha+\beta}{2}\right)\cdot\sin\left(\dfrac{\alpha-\beta}{2}\right)=\\
            =&2\cos\left(\dfrac{\alpha+\beta}{2}\right)\cdot\sin\left(\dfrac{\alpha-\beta}{2}\right)
        \end{align*}
    \end{proof}
    \begin{statement}
        $$\tg\alpha+\tg\beta=\dfrac{\sin(\alpha+\beta)}{\cos\alpha\cdot\cos\beta}$$
    \end{statement}
    \begin{proof}
        $$\tg\alpha+\tg\beta=\dfrac{\sin\alpha}{\cos\alpha}+\dfrac{\sin\beta}{\cos\beta}=\dfrac{\sin\alpha\cdot\cos\beta+\cos\alpha\cdot\sin\beta}{\cos\alpha\cdot\cos\beta}=\dfrac{\sin(\alpha+\beta)}{\cos\alpha\cdot\cos\beta}$$
    \end{proof}
    \begin{statement}
        $$\tg\alpha-\tg\beta=\dfrac{\sin(\alpha-\beta)}{\cos\alpha\cdot\cos\beta}$$
    \end{statement}
    \begin{proof}
        $$\tg\alpha-\tg\beta=\dfrac{\sin\alpha}{\cos\alpha}-\dfrac{\sin\beta}{\cos\beta}=\dfrac{\sin\alpha\cdot\cos\beta-\cos\alpha\cdot\sin\beta}{\cos\alpha\cdot\cos\beta}=\dfrac{\sin(\alpha-\beta)}{\cos\alpha\cdot\cos\beta}$$
    \end{proof}
    \begin{statement}
        $$\ctg\alpha+\ctg\beta=\dfrac{\sin(\alpha+\beta)}{\sin\alpha\cdot\sin\beta}$$
    \end{statement}
    \begin{proof}
        $$\ctg\alpha+\ctg\beta=\dfrac{\cos\alpha}{\sin\alpha}+\dfrac{\cos\beta}{\sin\beta}=\dfrac{\cos\alpha\cdot\sin\beta+\sin\alpha\cdot\cos\beta}{\sin\alpha\cdot\sin\beta}=\dfrac{\sin(\alpha+\beta)}{\sin\alpha\cdot\sin\beta}$$
    \end{proof}
    \begin{statement}
        $$\ctg\alpha-\ctg\beta=-\dfrac{\sin(\alpha-\beta)}{\sin\alpha\cdot\sin\beta}$$
    \end{statement}
    \begin{proof}
        $$\ctg\alpha-\ctg\beta=\dfrac{\cos\alpha}{\sin\alpha}-\dfrac{\cos\beta}{\sin\beta}=\dfrac{\cos\alpha\cdot\sin\beta-\sin\alpha\cdot\cos\beta}{\sin\alpha\cdot\sin\beta}=\dfrac{\sin(\beta-\alpha)}{\sin\alpha\cdot\sin\beta}$$
    \end{proof}
    \begin{statement}
        $$\cos\alpha\cdot\cos\beta=\dfrac{\cos(\alpha+\beta)+\cos(\alpha-\beta)}{2}$$
    \end{statement}
    \begin{proof}
        Пусть $\tilde{\alpha} = \dfrac{\alpha+\beta}{2}$, $\tilde{\beta} = \dfrac{\alpha-\beta}{2}$, тогда из утверждения 7.21:
        $$2\cos\tilde{\alpha}\cdot\cos\tilde{\beta}=\cos(\tilde{\alpha}+\tilde{\beta})+\cos(\tilde{\alpha}-\tilde{\beta})\Longleftrightarrow \cos\tilde{\alpha}\cdot\cos\tilde{\beta}=\dfrac{\cos(\tilde{\alpha}+\tilde{\beta})+\cos(\tilde{\alpha}-\tilde{\beta})}{2}$$
    \end{proof}
    \begin{statement}
        $$\sin\alpha\cdot\sin\beta=-\dfrac{\cos(\alpha+\beta)-\cos(\alpha-\beta)}{2}$$
    \end{statement}
    \begin{proof}
        Пусть $\tilde{\alpha} = \dfrac{\alpha+\beta}{2}$, $\tilde{\beta} = \dfrac{\alpha-\beta}{2}$, тогда из утверждения 7.22:
        $$-2\sin\tilde{\alpha}\cdot\sin\tilde{\beta}=\cos(\tilde{\alpha}+\tilde{\beta})-\cos(\tilde{\alpha}-\tilde{\beta})\Longleftrightarrow \sin\tilde{\alpha}\cdot\sin\tilde{\beta}=-\dfrac{\cos(\tilde{\alpha}+\tilde{\beta})-\cos(\tilde{\alpha}-\tilde{\beta})}{2}$$
    \end{proof}
    \begin{statement}
        $$\sin\alpha\cdot\cos\beta=\dfrac{\sin(\alpha+\beta)+\sin(\alpha-\beta)}{2}$$
    \end{statement}
    \begin{proof}
        Пусть $\tilde{\alpha} = \dfrac{\alpha+\beta}{2}$, $\tilde{\beta} = \dfrac{\alpha-\beta}{2}$, тогда из утверждения 7.23:
        $$2\sin\tilde{\alpha}\cdot\sin\tilde{\beta}=\sin(\tilde{\alpha}+\tilde{\beta})+\sin(\tilde{\alpha}-\tilde{\beta})\Longleftrightarrow \sin\tilde{\alpha}\cdot\cos\tilde{\beta}=\dfrac{\sin(\tilde{\alpha}+\tilde{\beta})+\sin(\tilde{\alpha}-\tilde{\beta})}{2}$$
    \end{proof}
    \begin{statement}
        $$\sin\alpha+\cos\alpha=\sqrt{2}\cdot\sin\left(\dfrac{\pi}{4}+\alpha\right)$$
    \end{statement}
    \begin{proof}
        $$\sin\alpha+\cos\alpha=\sqrt{2}\cdot\left(\dfrac{\sqrt{2}}{2}\cdot\sin\alpha+\dfrac{\sqrt{2}}{2}\cdot\cos\alpha\right)=\sqrt{2}\cdot\sin\left(\dfrac{\pi}{4}+\alpha\right)$$
    \end{proof}
    \begin{statement}
        $$\ctg\alpha=\dfrac{1+\cos(2\alpha)}{\sin(2\alpha)}$$
    \end{statement}
    \begin{proof}
        $$\ctg\alpha=\dfrac{\cos\alpha}{\sin\alpha}=\dfrac{2\cos^2\alpha}{2\sin\alpha\cdot\cos\alpha}=\dfrac{1+\cos(2\alpha)}{\sin(2\alpha)}$$
    \end{proof}
    \begin{statement}
        $$\tg\alpha=\dfrac{1-\cos(2\alpha)}{\sin(2\alpha)}$$
    \end{statement}
    \begin{proof}
        $$\tg\alpha=\dfrac{\sin\alpha}{\cos\alpha}=\dfrac{2\sin^2\alpha}{2\sin\alpha\cdot\cos\alpha}=\dfrac{1-\cos(2\alpha)}{\sin(2\alpha)}$$
    \end{proof}

    \section{Комплексные числа}
    \begin{definition}
        Мнимой единицей называется такое число $i$, что $i^2=-1$.
    \end{definition}
    \begin{definition}
        Комплексным числом называется выражение вида $z=a+bi$, где $a=\Re z$ – действительная часть, а $b=\Im z$ – мнимая.
    \end{definition}
    \begin{definition}
        Число $z_1$ называется сопряженным к числу $z$, если $\Re z_1=\Re z;\,\Im z_1=-\Im z$. Обозначение: $z_1=\bar{z}$.
    \end{definition}
    \begin{definition}
        Модулем $z\in \clx$ называется $|z|=\sqrt{(\Re z)^2+(\Im z)^2}$.
    \end{definition}
    \begin{definition}
        Комплексная плоскость:
        \begin{center}
            \begin{tikzpicture}[scale=1.2]
                \coordinate (O) at (0,0);
                \coordinate (a) at (3,0);
                \coordinate (b) at (0,2);
                \coordinate (z) at (3,2);
                \draw[-{Stealth[scale = 1.5]}] (-1.5,0) -- (5, 0) node[right] {$\Re z$};
                \draw[-{Stealth[scale = 1.5]}] (0,-1) -- (0, 4) node[above] {$\Im z$};
                \draw (O) -- (z);
                \draw[dashed] (a) -- (z) -- (b);
                \draw (O) node[below left] {$O$};
                \draw (a) node[below] {$a$};
                \draw (b) node[left] {$b$};
                \draw (z) node[above right] {$z$};
                \draw ($(O)!0.5!(z)$) node[above left] {$r$};
                \draw (3.5,3) node[above] {$z=r\cdot(\cos\varphi+ i\sin\varphi)$};
                \tkzLabelAngle(a,O,z){$\varphi$}
            \tkzMarkAngle[size=.55cm](a,O,z)
            \tkzDrawPoints(O, a, b, z)
        \end{tikzpicture}
        \end{center}
    \end{definition}
    \begin{theorem}[Формула Муавра]
        $$z^n=r^n\cdot(\cos(n\varphi)+i\sin(n\varphi))$$
    \end{theorem}
    \begin{proof}
        Докажем по индукции. $z=r\cdot(\cos\varphi+ i\sin\varphi)$ верно. Пусть $z^n=r^n\cdot(\cos(n\varphi)+i\sin(n\varphi))$ верно, докажем, что для $n+1$ утверждение верно верно.
        \begin{multline*}
            z^{n+1}=z^n\cdot z=r^n\cdot(\cos(n\varphi)+i\sin(n\varphi))\cdot r\cdot(\cos\varphi+ i\sin\varphi)=\\
            =r^{n+1}\cdot((\cos(n\varphi)\cdot\cos\varphi-\sin(n\varphi)\cdot\sin\varphi)+i(\cos(n\varphi)\cdot\sin\varphi+\sin(n\varphi)\cdot\cos\varphi))=\\
            =r^{n+1}\cdot(\cos((n+1)\varphi)+i\sin((n+1)\varphi))
        \end{multline*}
    \end{proof}
    \begin{definition}
        Корнем из комплексного числа называется:
        $$\sqrt[n]{z}=\sqrt[n]{r}\left(\cos\left(\frac{\varphi}{n}+\frac{2\pi k}{n}\right)+i\sin\left(\frac{\varphi}{n}+\frac{2\pi k}{n}\right)\right),\,k\in\{1,\,2,\ldots,\,n-1\}$$
    \end{definition}
    \subsection{Матрицы поворота}
    \begin{definition}
        Матрицей поворота на угол $\alpha$ называется:
        $$R_\alpha=\begin{pmatrix}
            \cos\alpha & -\sin\alpha\\
            \sin\alpha & \cos\alpha
        \end{pmatrix}
        =\cos \alpha \cdot \begin{pmatrix}
            1 & 0\\
            0 & 1
        \end{pmatrix} +\sin \alpha \cdot \begin{pmatrix}
            0 & -1\\
            1 & 0
        \end{pmatrix}=\cos \alpha\cdot E+\sin\alpha\cdot I$$
    \end{definition}
    \begin{statement}
        $$z=r\cdot(\cos\varphi+ i\sin\varphi)\Longleftrightarrow z=re^{i\varphi}$$
    \end{statement}
    \begin{proof}
        $$e^x=\lim_{n\to\infty}\sum_{i=0}^n\dfrac{x^i}{i!}$$
        $$\cos x = \lim_{n\to\infty}\sum_{i=0}^n\dfrac{x^{2i}}{(2i)!}\cdot(-1)^i$$
        $$\sin x = \lim_{n\to\infty}\sum_{i=0}^n\dfrac{x^{2i+1}}{(2i+1)!}\cdot(-1)^i$$
        $$e^{ix}=\lim_{n\to\infty}\sum_{k=0}^n\dfrac{x^k}{k!}\cdot i^k=\cos x+ i\sin x$$
    \end{proof}

    \section{Логарифмы}
    \begin{definition}
        Пусть $a>0;\, a\neq 1,\,b>0$. Логарифмом $\log_ab$ числа $b$ по основанию $a$ называется такое число $c$, что $a^c=b$. Из этого, а также свойств степеней следует:
        \begin{center}
            \begin{multicols}{3}
                \[\log_a 1=0\]\\
                \[\log_a a=1\]\\
                \[a^{\log_a b}=b\]
            \end{multicols}
        \end{center}
    \end{definition}
    \begin{statement}
        $$\log_a b+\log_a c=\log_a(bc)$$
    \end{statement}
    \begin{proof}
        $$a^{\log_a b+\log_a c}=a^{\log_a b}\cdot a^{\log_a c}=bc=a^{\log_a(bc)}$$
    \end{proof}
    \begin{statement}
        $$\log_a b-\log_a c=\log_a \left(\frac{b}{c}\right)$$
    \end{statement}
    \begin{proof}
        $$a^{\log_a b-\log_a c}=\frac{a^{\log_a b}}{a^{\log_a c}}=\frac{b}{c}=a^{\log_a \left(\frac{b}{c}\right)}$$
    \end{proof}
    \begin{statement}
        $$\log_a \left(b^k\right)=k\cdot\log_a b$$
    \end{statement}
    \begin{proof}
        $$a^{\log_a \left(b^k\right)}=b^k=\left(a^{\log_a b}\right)^k=a^{k\cdot\log_a b}$$
    \end{proof}
    \begin{statement}
        $$\log_{a^k}b=\frac{1}{k}\cdot\log_a b$$
    \end{statement}
    \begin{proof}
        $$\left(a^k\right)^{\log_{a^k}b}=a^{k\cdot\log_{a^k} b}=b=a^{\log_a b}$$
    \end{proof}
    \setcounter{subsection}{4}
    \begin{consequence}
        $$\log_{a^k}\left(b^k\right)=\log_a b$$
    \end{consequence}
    \begin{statement}
        $$\log_ab=\frac{\log_cb}{\log_ca},\,\log_ca\neq0$$
    \end{statement}
    \begin{proof}
        $$\log_ab=\frac{\log_cb}{\log_ca}\Longleftrightarrow\log_ab\cdot\log_ca=\log_cb$$
        $$c^{\log_cb}=b=a^{\log_ab}=\left(c^{\log_ca}\right)^{\log_ab}=c^{\log_ab\cdot\log_ca}$$
    \end{proof}
    \begin{statement}
        $$\log_ab=\frac{1}{\log_ba}$$
    \end{statement}
    \begin{proof}
        $$a^{\frac{1}{\log_ba}}=\left(b^{\log_ba}\right)^{\frac{1}{\log_ba}}=b=a^{\log_ab}$$
    \end{proof}
    \begin{statement}
        $$a^{\log_bc}=c^{\log_ba}$$
    \end{statement}
    \begin{proof}
        $$a^{\log_bc}=a^{\frac{\log_ac}{\log_ab}}=\left(a^{\log_ac}\right)^{\frac{1}{\log_ab}}=c^\frac{1}{\log_ab}=c^{\log_ba}$$
    \end{proof}
    \begin{definition}
        Натуральным логарифмом $\ln x$ называется логарифм с основанием $e$.
    \end{definition}
    \begin{definition}
        Десятичным логарифмом $\lg x$ называется логарифм с основанием 10.
    \end{definition}

    \subsection{Показательные уравнения и неравенства}

    \begin{tasks}(2)
    \task[] $a^x=a^{x_0}\Longleftrightarrow x=x_0$\\
    \task[] $a^x=b\Longleftrightarrow x=\log_ab$\\
    \task[] $a^x>a^{x_0}\Longleftrightarrow\begin{array}{l}
        x>x_0\text{, если }a>1\\
        x<x_0\text{, если }0<a<1
    \end{array}$
    \task[] $a^x>b\Longleftrightarrow\begin{array}{l}
        x>\log_ab\text{, если }a>1\\
        x<\log_ab\text{, если }0<a<1
    \end{array}$
    \end{tasks}

    \subsection{Логарифмические уравнения и неравенства}

    \begin{tasks}(2)
        \task[] $\log_ax=b\Longleftrightarrow x=a^b$\\
        \task[] $\log_ax>b\Longleftrightarrow \begin{array}{l}
            x>a^b\text{, если }a>1\\
            x<a^b\text{, если }0<a<1
        \end{array}$\\
    \end{tasks}
    $$\log_af(x)<\log_ag(x)\Longleftrightarrow
        \begin{array}{l}
            \begin{cases}
                f(x) < g(x)\\
                f(x) > 0
            \end{cases}a>1\\
            \begin{cases}
                f(x) > g(x)\\
                g(x) > 0
            \end{cases}0<a<1\\
        \end{array}$$

\section{Линейная алгебра}

\subsection{Координаты на плоскости}

\begin{statement}
    Пусть даны $A(x_1;y_1)$ и $B(x_2;y_2)$. Тогда середина отрезка $AB$ будет иметь координаты $(\frac{x_1+x_2}{2};\frac{y_1+y_2}{2})$, а его длина будет равна $d=\sqrt{(x_2-x_1)^2+(y_2-y_1)^2}$.
\end{statement}
\begin{statement}
    Пусть $A(x_0;y_0)\in\omega(O;r),\,O(a;b)$, тогда $(x_0-a)^2+(y_0-b)^2=r^2$.
\end{statement}
\begin{statement}
    Пусть даны $A(x_1;y_1)$ и $B(x_2;y_2)$. Тогда каноническое уравнение прямой, проходящей через $A$ и $B$ будет иметь вид:
    $$\frac{x-x_1}{x_2-x_1}=\frac{y-y_1}{y_2-y_1}$$
\end{statement}
\begin{statement}
    Пусть даны точка $M(x_0;y_0)$ и прямая $l:\, ax+by+c=0$. Тогда:
    $$\rho(M;l)=\frac{|ax_0+by_0+c|}{\sqrt{a^2+b^2}}$$
\end{statement}

\subsection{Векторы}

\begin{definition}
    Направленным отрезком $\vv{AB}$ называется отрезок $AB$, у которого заданы начало $A$ и конец $B$.
\end{definition}
\begin{definition}
    Направленные отрезки $\vv{AB}$ и $\vv{CD}$ называются коллинеарными, если $AB\parallel CD$.
\end{definition}
\begin{definition}
    Направленные отрезки $\vv{AB}$ и $\vv{CD}$ называются конгруэнтными, если они коллинеарны, $B$ и $D$ лежат в одной плоскости относительно $AC$ и их длины равны.
\end{definition}
\begin{definition}
    Пусть $\vv{r_1},\ldots,\vv{r_n}$ – векторы в $\R^n$. Они называются линейно зависимыми, если $\exists\,a_1,\ldots,a_n$, не равные одновременно нулю, что $a_1\vv{r_1}+\ldots+a_n\vv{r_n}=0$.
\end{definition}
\begin{definition}
    Пусть $X$ – множество, рассмотрим $X\times X=\{(x;y)\mid x,y\in X\}$. Тогда $\rho\subset X\times X$ – бинарное отношение.
\end{definition}
\begin{definition}
    Бинарное отношение $\sim\subset X\times X$ называется отношением эквивалентности на множестве $X$, если $\forall\,x,y,z\in X:\,x\sim x;\,x\sim y\iff y\sim x;\,x\sim y,\,y\sim z\iff x\sim z$.
\end{definition}

\subsection{Линейные пространства}

\begin{definition}
    Элементами поля $F$ являются скаляры или векторы. $L$ называется линейным пространством над полем $F$, если $\forall\,a,b\in L;\,\forall\,\lambda,\,\mu \in F$:
    \begin{center}
        \begin{multicols}{2}
            \begin{enumerate}
                \item $a+b=b+a$\\
                \item $(a+b)+c=a+(b+c)$\\
                \item $\exists\,\vv{0}\in L:\,a+\vv{0}=a$\\
                \item $\forall\,a\enspace  \exists\,b:\,a+b=\vv{0}$\\
                \item $\lambda(\mu a)=(\lambda \mu)a$\\
                \item $(\lambda+\mu)a=\lambda a + \mu a$\\
                \item $\lambda(a+b)=\lambda a+ \lambda b$\\
                \item $1\cdot a = a$
            \end{enumerate}
        \end{multicols}
    \end{center}
\end{definition}
\begin{definition}
    Пусть $L$ – линейное пространство над полем $F$; $u_1,\ldots,u_n\in L;\,\alpha_1,\ldots,\alpha_n\in F$, тогда $\alpha_1u_1+\ldots+\alpha_nu_n$ – линейная комбинация. Линейная комбинация называется тривиальной, если $\forall\,i:\,\alpha_i=0$.
\end{definition}
\begin{definition}
    Пусть $U\subset L$, множество векторов $\langle U\rangle$, которые не выражаются через элементы $U$ называется линейной оболочкой.
\end{definition}
\begin{definition}
    Пусть $\{x_1,\ldots,x_n\}$ – система векторов. Она называется линейно зависимой, если существует нетривиальная линейная комбинация, равная нулю.
\end{definition}
\begin{theorem}
    Система векторов является линейно зависимой тогда и только тогда, когда один из векторов является линейной комбинацией других.
\end{theorem}
\begin{definition}
    Система $U=\{u_1,\ldots,u_n\}$ линейно независимых векторов линейного пространства $L$ называется базисом, если $\langle U \rangle=L$.
\end{definition}
\begin{definition}
    Скалярным произведением векторов $\vv{a}\neq\vv{0}$ и $\vv{b}\neq\vv{0}$ называется величина $|\vv{a}|\cdot|\vv{b}|\cdot\cos\angle(\vv{a};\vv{b})$. Если $\vv{a}$ и $\vv{b}$ сонаправлены, то угол между ними равен нулю, а если коллинеарны, но не сонаправлены – $\pi$:
    \begin{center}
        \begin{multicols}{2}
            \begin{enumerate}
                \item $\vv{a}\cdot\vv{b}=\vv{b}\cdot\vv{a}$\\
                \item $k\cdot(\vv{a};\vv{b})=(k\cdot\vv{a})\cdot\vv{b}=\vv{a}\cdot(k\cdot\vv{b})$\\
                \item $\vv{a}\cdot(\vv{b}+\vv{c})=\vv{a}\cdot\vv{b}+\vv{a}\cdot\vv{c}$\\
                \item $\vv{a}\cdot\vv{b}=\dfrac{1}{2}\left(|\vv{a}+\vv{b}|^2-|\vv{a}|^2-|\vv{b}|^2\right)$
            \end{enumerate}
        \end{multicols}
    \end{center}
\end{definition}

\subsection{Матрицы}

\begin{definition}
    Квадратной матрицей называется матрица, у которой число строк равно числу столбцов.
\end{definition}
\begin{definition}
    Квадратная матрица называется верхней треугольной, если $\forall\,i,j;\,i>j:\,a_{ij}=0$.
\end{definition}
\begin{definition}
    Квадратная матрица называется нижней треугольной, если $\forall\,i,j;\,i<j:\,a_{ij}=0$.
\end{definition}
\begin{definition}
    Квадратная матрица называется диагональной, если $\forall\,i,j;\,i\neq j:\,a_{ij}=0$. Обозначение: $\diag(a_{11},\ldots,a_{nn})$.
\end{definition}
\begin{definition}
    Диагональная матрица порядка $n$, у которой все диагональные элементы равны 1, называется единичной. Обозначение: $E$ или $E_n$. Элементы единичной матрицы обозначаются $\delta_{ij}$:
    $$\delta_{ij}=\left\{\begin{alignedat}{2}
        & 1 && \text{ при }i=j \\
        & 0  && \text{ при }i\neq j
      \end{alignedat}\right.$$
\end{definition}
\begin{definition}
    Квадратная матрица $A$ называется скалярной, если $A=\diag(\lambda,\ldots,\lambda)$.
\end{definition}
\begin{definition}
    Сумма диагональных элементов матрицы $A$ называется её шпуром или следом. Обозначение: $\spur A$ или $\tr A$.
\end{definition}
\begin{definition}
    Матрица $B$ называется транспонированной по отношению к матрице $A$, если $\forall\,i,j:\,b_{ij}=a_{ji}$. Обозначение: $B=A^T$.
\end{definition}
\begin{definition}
    Квадратная матрица $A$ называется симметрической, если $A^T=A$.
\end{definition}
\begin{definition}
    Квадратная матрица $A$ называется кососимметрической, если $A^T=-A$.
\end{definition}
\begin{definition}
    Матрица $B$ называется комплексно сопряжённой по отношению к матрице $A$, если $\forall\,i,j:\,b_{ij}=\overline{a_{ij}}$. Обозначение: $B=\overline{A}$.
\end{definition}
\begin{definition}
    Матрица $B$ называется эрмитово сопряжённой по отношению к матрице $A$, если $\forall\,i,j:\,b_{ij}=\overline{a_{ji}}$. Обозначение: $B=A^*$.
\end{definition}
\begin{definition}
    Квадратная матрица $A$ называется эрмитовой, если $A^*=A$.
\end{definition}
\begin{definition}
    Квадратная матрица $A$ называется косоэрмитовой, если $A^*=-A$.
\end{definition}
\begin{definition}
    Матрица называется нулевой, если все её элементы равны нулю. Обозначение: $A=O$.
\end{definition}
\begin{definition}
    Матрица называется неотрицательной, если $\forall\,i,j:\,a_{ij}\geq0$.
\end{definition}
\begin{definition}
    Матрица называется стохастической, если:
    $$\forall\,i,j:\,a_{ij}\geq0;\,\forall\,i\in\{1,\ldots,n\}:\,\sum_{k=1}^{n}a_{ik}=1$$
\end{definition}

\subsubsection{Произведение матриц}

\begin{definition}
    Пусть даны матрица $A$ размера $m\times n$ и матрица $B$ размера $n\times k$. Их произведением будет матрица размера $m\times k$, в которой:
    $$c_{ij}=\sum_{l=1}^{n}a_{il}\cdot b_{lj}$$
\end{definition}
\begin{statement}
    Произведение верхних треугольных матриц является верхней треугольной матрицей.
\end{statement}
\begin{definition}
    Пусть $A$ – квадратная матрица порядка $n$. Матрица $B$ называется обратной к $A$, если $A\cdot B=B\cdot A=E$. Обозначение: $B=A^{-1}$.
\end{definition}
\begin{definition}
    Квадратная матрица $A$ называется ортогональной, если $A^T=A^{-1}$.
\end{definition}
\begin{definition}
    Квадратная матрица $A$ называется унитарной, если $A^*=A^{-1}$.
\end{definition}
\begin{definition}
    Матрица $A$ называется нильпотентной, если $A^k=O,\,k\in \N$. Наименьшее из таких $k$ называется показателем нильпотентности матрицы $A$.
\end{definition}
\begin{definition}
    Матрица $A$ называется периодической, если $A^k=E,\,k\in \N$. Наименьшее из таких $k$ называется периодом матрицы $A$.
\end{definition}

\subsubsection{Определитель}

\begin{definition}
    Пусть дана перестановка $\begin{pmatrix}
        1 & 2 & \cdots & n\\
        \alpha_1 & \alpha_2 & \cdots & \alpha_n
    \end{pmatrix}$. Инверсией называется число таких пар $(i;j)$, что $i>j;\,\alpha_i<\alpha_j$.
\end{definition}
\begin{definition}
    Пусть дана квадратная матрица $A$ порядка $n$. Если $\alpha_i$ – перестановка чисел от 1 до $n$, а $N(\alpha_i)$ – число инверсия в $\alpha_i$ перестановке, то определителем матрицы $A$ называется:
    $$\det A=\sum_{a_i}(-1)^{N(\alpha_i)}\cdot a_{1\alpha_1}\cdot a_{2\alpha_2}\cdot\ldots\cdot a_{n\alpha_n}$$
\end{definition}
\begin{theorem}
    Пусть дана квадратная матрица $A$, $M_{ji}$ – дополнительный минор элемента $a_{ji}$. Элементы её обратной матрицы можно вычислить по формуле:
    $$b_{ij}=\frac{(-1)^{i+j}\cdot M_{ji}}{\det A}$$
\end{theorem}
\begin{definition}
    Квадратная матрица $A$ называется вырожденной, если $\det A=0$.
\end{definition}
\begin{definition}
    Квадратная матрица $A$ называется невырожденной, если $\det A\neq0$.
\end{definition}
\begin{definition}
    Квадратная матрица $A$ называется унимодулярной, если $|\det A|=1,\,a_{ij}\in\Z$.
\end{definition}
\begin{definition}
    Квадратная матрица $A$ называется матрицей перестановки, если она получена из $E$ перестановкой строк.
\end{definition}
\begin{definition}
    Квадратная матрица $A$ называется элементарной, если она получена из $E$ элементарным преобразованием.
\end{definition}

\subsubsection{Ранг матрицы}

\begin{definition}
    Рангом матрицы называется наибольший порядок из всех порядков её ненулевых миноров.
\end{definition}
\begin{theorem}[О базисном миноре]
    Пусть $M_r$ – базисный минор матрицы $A$, строки соответствующей матрицы – базисные строки, а столбцы – базисные столбцы. Тогда:
    \begin{multicols}{2}
        \begin{enumerate}
        \item \label{базисный минор 1} Базисные строки линейно независимы.
        \item \label{базисный минор 2} Базисные столбцы линейно независимы.
        \item Любая строка (столбец) матрицы $A$ – линейная комбинация базисных строк (столбцов).
    \end{enumerate}
    \end{multicols}
\end{theorem}
\begin{proof}
    Пункты \ref{базисный минор 1} и \ref{базисный минор 2} следуют из определения базиса. Без ограничения общности пусть базисный минор находится в левом верхнем углу матрицы $A$, то есть матрица имеет вид:
    $$\begin{pmatrix}
        \boldsymbol{a_{11}} & \boldsymbol{a_{12}} & \cdots & \boldsymbol{a_{1r}} & a_{1(r+1)} & \cdots & a_{1n}\\
        \boldsymbol{a_{21}} & \boldsymbol{a_{22}} & \cdots & \boldsymbol{a_{2r}} & a_{2(r+1)} & \cdots & a_{2n}\\
        \vdots & \vdots & \ddots & \vdots & \vdots & \ddots & \vdots\\
        \boldsymbol{a_{r1}} & \boldsymbol{a_{r2}} & \cdots & \boldsymbol{a_{rr}} & a_{r(r+1)} & \cdots & a_{rn}\\
        a_{(r+1)1} & a_{(r+1)2} & \cdots & a_{(r+1)r} & a_{(r+1)(r+1)} & \cdots & a_{(r+1)n}\\
        \vdots & \vdots & \ddots & \vdots & \vdots & \ddots & \vdots\\
        a_{m1} & a_{m2} & \cdots & a_{mr} & a_{m(r+1)} & \cdots & a_{mn}
    \end{pmatrix}$$
    Если расширить базисный минор до $r+1$ строки и $r+1$ столбца, при этом все его строки и столбцы останутся линейно независимыми, тогда $\rank A=r+1$, но по условию $\rank A=r$, противоречие.
\end{proof}
\begin{theorem}[Кронекер-Капелли]
    Система линейных алгебраических уравнений $A\cdot \vv{x}=\vv{b}$ совместна тогда и только тогда, когда $\rank A=\rank\left(A\mid \vv{b}\right)$.
\end{theorem}
\begin{proof}
    Пусть система совместна, тогда:
    $$b_i=\sum_{j}a_{ji}x_j;\, \vv{b}=\sum_{j}\vv{a_j}x_j$$
    То есть, $\vv{b}$ выражается через $\vv{a_i}$, а значит $\rank A=\rank\left(A\mid \vv{b}\right)$.
    Пусть $\rank A=\rank\left(A\mid \vv{b}\right)$. Тогда по теореме о базисном миноре она совместна.
\end{proof}
\begin{theorem}[Правило Крамера]
    Пусть дана система из $n$ линейных уравнений с $n$ неизвестным вида $A\cdot \vv{x}=\vv{b}$:
    $$\begin{cases}
        a_{11}x_1 + a_{12}x_2 + \ldots + a_{1n}x_n = b_1\\
        a_{21}x_1 + a_{22}x_2 + \ldots + a_{2n}x_n = b_2\\
        \makebox[57mm]{\dotfill}\\
        a_{n1}x_1 + a_{n2}x_2 + \ldots + a_{nn}x_n = b_n
    \end{cases}$$
    Тогда если $\Delta$ – определитель матрицы $A$, то $x_i=\dfrac{\Delta_i}{\Delta}$, где $\Delta_i$ – определитель матрицы $A$, в которой $i$-й столбец заменили на вектор свободных членов:
    $$x_i=\frac{1}{\Delta}\begin{vmatrix}
        a_{11} & \cdots & a_{1(i-1)} & b_1 & a_{1(i+1)} & \cdots & a_{1n}\\
        a_{21} & \cdots & a_{2(i-1)} & b_2 & a_{2(i+1)} & \cdots & a_{2n}\\
        \vdots & \ddots & \vdots & \vdots & \vdots & \ddots & \vdots\\
        a_{(n-1)1} & \cdots & a_{(n-1)(i-1)} & b_{(n-1)} & a_{(n-1)(i+1)} & \cdots & a_{(n-1)n}\\
        a_{n1} & \cdots & a_{n(i-1)} & b_n & a_{n(i+1)} & \cdots & a_{nn}
    \end{vmatrix}$$
    Если $\Delta=0$, то система имеет бесконечно много решений или несовместна.
\end{theorem}
\begin{theorem}[Конечномерная альтернатива Фредгольма]
    Пусть дана система из $m$ линейных уравнений с $n$ неизвестными вида $A\cdot \vv{x}=\vv{b}$:
    $$\begin{cases}
        a_{11}x_1 + a_{12}x_2 + \ldots + a_{1n}x_n = b_1\\
        a_{21}x_1 + a_{22}x_2 + \ldots + a_{2n}x_n = b_2\\
        \makebox[57mm]{\dotfill}\\
        a_{m1}x_1 + a_{m2}x_2 + \ldots + a_{mn}x_n = b_m
    \end{cases}$$
    Тогда выполняется одно из двух условий:
    \begin{center}
        \begin{multicols}{2}
            \begin{enumerate}
                \item Система имеет 1 решение для любого $\vv{b}$ и соответствующая однородная система уравнений имеет только тривиальное решение.
                \item Соответствующая однородная система уравнений имеет нетривиальное решение, тогда $\exists\,\vv{b}:\,A\cdot\vv{x}=\vv{b}$ не имеет решений.
            \end{enumerate}
        \end{multicols}
    \end{center}
\end{theorem}
\begin{proof}
    \hfill
    \begin{enumerate}
        \item Пусть $\det A\neq0$. Тогда все столбцы $A$ называются векторами $n$-мерного пространства, при этом они линейно независимы, то есть они образуют в этом пространстве базис. Тогда $\vv{b}=x_1\vv{a_1}+x_2\vv{a_2}+\ldots+x_n\vv{a_n}$; любой вектор этого пространства выражается единственным образом через линейную комбинацию базисных векторов. Система имеет единственное решение. Поскольку векторы $\vv{a_1},\vv{a_2},\ldots,\vv{a_n}$ линейно независимы, то нулю может равняться только их тривиальная комбинация.
        \item Пусть $\det A=0$. Векторы $\vv{a_1},\vv{a_2},\ldots,\vv{a_n}$ – линейно зависимы, поэтому существует их нетривиальная линейная комбинация, равная нулю. $\dim \langle \vv{a_i} \rangle=k<n$. Пусть при любом $\vv{b}$ система имеет решение. Тогда любой $\vv{b}$ выражается через линейную комбинацию $\vv{a_i}$, то есть $\{\vv{a_i}\}$ – базис в $n$-мерном пространстве. Противоречие.
    \end{enumerate}
\end{proof}

\end{document}